\makeatletter
\def\input@path{{styles/}}
\makeatother
\ifx\bookMode\undefined%
    %-----------------------------------------------------------------------------
% Master Prefix File for Randomized Algorithms Class Notes & Book
% Unified self-contained configuration for LuaLaTeX and XeLaTeX with Memoir
%-----------------------------------------------------------------------------

\RequirePackage{iftex}
\ifPDFTeX
  \PackageError{engine-guard}{This document requires LuaLaTeX or XeLaTeX}{%
    Please recompile using LuaLaTeX or XeLaTeX instead of pdfLaTeX.%
  }
  \stop
\fi

% Always use memoir
\documentclass[12pt,oneside,fleqn]{memoir}

%-----------------------------------------------------------------------------
% Typography, Fonts & Math Packages
%-----------------------------------------------------------------------------
\usepackage{mathtools}
\usepackage{newcomputermodern}   % Loads fontspec & unicode-math

% OpenType math aliases and fixes
\providecommand{\square}{\mdlgwhtsquare}
\def\Box{\mdlgwhtsquare}
\def\Diamond{\mdlgwhtdiamond}
\def\blacksquare{\mdlgblksquare}
\def\leadsto{\rightsquigarrow}

\usepackage{alphabeta}
\usepackage[normalem]{ulem} % Provides \sout
\usepackage{mleftright}
\usepackage{microtype}
\usepackage{xspace}
\usepackage{xurl}
\usepackage[dvipsnames,table]{xcolor}
\usepackage{graphicx}
\usepackage{caption}
\usepackage{booktabs,array,colortbl,hhline,multirow,tabularx}
\usepackage[inline]{enumitem}
\usepackage[thicklines]{cancel}
\renewcommand\CancelColor{\color{purple}}
\usepackage{euscript}
\usepackage{marvosym}
\let\Cross\undefined
\let\Square\undefined
\usepackage{bbding}
\usepackage{stmaryrd}
\usepackage{epigraph}
\usepackage{makeidx}
\usepackage{perpage}
\usepackage{tikz}

%-----------------------------------------------------------------------------
% Boondox / mathcalb font declaration
%-----------------------------------------------------------------------------
\DeclareFontFamily{U}{BOONDOX-calo}{\skewchar\font=45 }
\DeclareFontShape{U}{BOONDOX-calo}{m}{n}{
  <-> s*[1.05] BOONDOX-r-calo}{}
\DeclareFontShape{U}{BOONDOX-calo}{b}{n}{
  <-> s*[1.05] BOONDOX-b-calo}{}
\DeclareMathAlphabet{\mathcalb}{U}{BOONDOX-calo}{m}{n}
\SetMathAlphabet{\mathcalb}{bold}{U}{BOONDOX-calo}{b}{n}
\DeclareMathAlphabet{\mathbcalb}{U}{BOONDOX-calo}{b}{n}

%-----------------------------------------------------------------------------
% Algorithm Environment (formerly salgorithm.sty)
%-----------------------------------------------------------------------------
\newcommand{\ProgWordX}[1]{{\small\textbf{#1}}{\ }}
\newcommand{\Do}{\ProgWordX{do}\xspace}
\newcommand{\End}{\ProgWordX{end}\xspace}
\newcommand{\Return}{\ProgWordX{return}\xspace}%
\newcommand{\Procbegin}{\ProgWordX{begin}}
\newcommand{\If}{\ProgWordX{if}\xspace}
\newcommand{\Then}{\ProgWordX{then}\xspace}
\newcommand{\Else}{\ProgWordX{else}\xspace}
\newcommand{\For}{\ProgWordX{for}\xspace}
\newcommand{\While}{{\small\bfseries while}\ }

\newcommand{\xbeginlgox}{\begin{minipage}{1in}\begin{tabbing}
           \quad\=\qquad\=\qquad\=\qquad\=\qquad\=\qquad\=\qquad\=\kill}
\newcommand{\xendlgox}{\end{tabbing}\end{minipage}}

\newenvironment{algorithmEnv}{%
   \begin{tabular}{|l|}\hline\xbeginlgox\bigskip}%
   {\xendlgox\\\hline\end{tabular}}

\newenvironment{program}{%
   \begin{minipage}{4.0in}%
   \begin{tabbing}%
       \ \ \ \ \=\ \ \ \=\ \ \ \ \=\ \ \ \ \=\ \ \ \ \=%
      \ \ \ \ \=\ \ \ \ \=\ \ \ \ \=\ \ \ \ \=%
      \ \ \ \ \=\ \ \ \ \=\ \ \ \ \=\ \ \ \ \=\kill%
}{%
   \end{tabbing}%
   \end{minipage}%
}

\newcommand{\AlgorithmBox}[1]{%
   \fbox{\begin{program} #1 \end{program}}%
}

\definecolor{AlgorithmColor}{cmyk}{0.07,0.90,0,0.34}
\providecommand{\Algorithm}[1]{{\AlgorithmI{#1}\index{algorithm!#1@{\AlgorithmI{#1}}}}}
\newcommand{\CodeComment}[1]{\textcolor{blue}{\texttt{#1}}}

%-----------------------------------------------------------------------------
% Path & Directory Resolution
%-----------------------------------------------------------------------------
\newcommand{\RootDIR}{../}
\IfFileExists{../../../book.tex}{\renewcommand{\RootDIR}{../../../}}{}
\IfFileExists{../../book.tex}{\renewcommand{\RootDIR}{../../}}{}
\IfFileExists{../book.tex}{\renewcommand{\RootDIR}{../}}{}
\IfFileExists{./book.tex}{\renewcommand{\RootDIR}{}}{}

\providecommand{\File}[1]{#1}
\providecommand{\RootFile}[1]{\RootDIR #1}

%-----------------------------------------------------------------------------
% QED Duck Symbol
%-----------------------------------------------------------------------------
\newlength{\CharHeight}
\AtBeginDocument{\setlength{\CharHeight}{\fontcharht\font`X}}

\IfFileExists{\RootDIR shared/qed_duck.pdf}{%
   \providecommand{\myqedsymbol}{%
      \reflectbox{\includegraphics[height=1.4\CharHeight]{\RootDIR shared/qed_duck.pdf}}%
   }%
}{%
   \IfFileExists{shared/qed_duck.pdf}{%
      \providecommand{\myqedsymbol}{%
         \reflectbox{\includegraphics[height=1.4\CharHeight]{shared/qed_duck.pdf}}%
      }%
   }{%
      \providecommand{\myqedsymbol}{$\blacksquare$}%
   }%
}

%-----------------------------------------------------------------------------
% BibLaTeX Setup (Pure Biber, No BibTeX)
%-----------------------------------------------------------------------------
\usepackage[bibencoding=utf8,style=alphabetic,refsegment=chapter,
indexing=cite,
backend=biber]{biblatex}

\ExecuteBibliographyOptions{%
   maxnames = 40,%
   maxalphanames = 4,%
   maxcitenames = 3,%
   minalphanames = 3,%
   minnames = 1,%
   giveninits = true,%
   doi = false,%
   url = false,%
   isbn = false%
}

\renewbibmacro{in:}{}

\AtEveryBibitem{%
   \ifentrytype{inproceedings}{%
      \clearname{editor}%
      \clearlist{editor}%
      \clearfield{editor}%
      \clearlist{publisher}%
      \clearfield{series}%
      \clearlist{location}%
      \clearfield{publisher}%
   }{}%
   \ifentrytype{book}{%
      \clearname{publisher}%
   }{}%
}
\AtEveryCitekey{%
   \ifentrytype{inproceedings}{%
      \clearlist{publisher}%
   }{}%
   \clearname{editor}%
   \clearlist{location}%
}
\DeclareFieldFormat[inproceedings]{date}{}
\DeclareFieldFormat[inproceedings]{pages}{#1, \printfield{year}}
\DeclareFieldFormat[inproceedings]{booktitle}{\mkbibemph{#1,\nopunct}}

\DeclareFieldFormat{journaltitle}{\mkbibemph{#1}\isdot}

\DeclareFieldFormat[article]{number}{\mkbibparens{#1}}
\DeclareFieldFormat[article]{date}{}
\DeclareFieldFormat[article]{pages}{#1, \printfield{year}}
\DeclareFieldFormat[article]{issue}{\nopunct#1}
\renewbibmacro*{issue+date}{%
   \iffieldundef{issue}{}%
   {%
      \printtext{\unspace}%
      (\printfield{issue})%
   }%
   \iffieldundef{pages}%
   {\addcomma\printfield{year}}%
   {}%
}

\renewcommand{\bibpagespunct}{%
   \ifentrytype{article}{\addcolon}{\addcomma}\space%
}

\DeclareFieldFormat[article]{titlecase}{\MakeSentenceCase{#1}}
\DeclareFieldFormat[inproceedings]{titlecase}{\MakeSentenceCase{#1}}
\DeclareFieldFormat[incollection]{titlecase}{\MakeSentenceCase{#1}}
\DeclareFieldFormat[book]{titlecase}{\MakeSentenceCase{#1}}

\renewbibmacro*{journal}{%
   \iffieldundef{journaltitle} {} {\printtext[journaltitle]{%
         \printfield[myplain]{journaltitle}%
         \setunit{\subtitlepunct}%
         \printfield[myplain]{journalsubtitle}}}}
\DeclareFieldFormat{myplain}{#1}

\renewbibmacro*{booktitle}{%
   \ifboolexpr{ test {\iffieldundef{booktitle}} and test
      {\iffieldundef{booksubtitle}} } {} {\printtext[booktitle]{%
         \printfield[myplain]{booktitle}%
         \setunit{\subtitlepunct}%
         \printfield[myplain]{booksubtitle}}%
      \newunit}%
   \printfield{booktitleaddon}}

\renewbibmacro*{volume+number+eid}{%
   \setunit*{\addcomma\addspace}%
   \printfield{volume}%
   \printfield{number}%
   \setunit{\addcomma\space}%
   \printfield{eid}%
}

\newbibmacro{string+doi}[1]{%
   \iffieldundef{doi}{%
      \iffieldundef{url}{%
         #1%
      }{%
         \href{\thefield{url}}{#1}%
      }%
   }{%
      \href{http://dx.doi.org/\thefield{doi}}{#1}}%
}
\DeclareFieldFormat*{title}{\usebibmacro{string+doi}{\mkbibemph{#1}}}
\renewcommand*{\multicitedelim}{\addcomma\space}
\DeclareFieldFormat[inproceedings]{title}{\usebibmacro{string+doi}{{#1}}}
\DeclareFieldFormat[article]{title}{\usebibmacro{string+doi}{{#1}}}
\DeclareFieldFormat[incollection]{title}{\usebibmacro{string+doi}{{#1}}}
\DeclareFieldFormat[book]{titlecase}{{#1}}

\DeclareCiteCommand{\tcite}%
{\usebibmacro{prenote}}%
{\usebibmacro{citeindex}\usebibmacro{cite}}%
{\multicitedelim}%
{\usebibmacro{postnote}}

%-----------------------------------------------------------------------------
% Hyperref & Visual Formatting
%-----------------------------------------------------------------------------
\usepackage{hyperref}
\hypersetup{%
   breaklinks,%
   colorlinks=true,%
   linkcolor=[rgb]{0.45,0.0,0.0},%
   citecolor=[rgb]{0,0,0.45}%
}

\definecolor{DeepNavy}{HTML}{002366}
\definecolor{RoyalBlue}{HTML}{004E92}
\definecolor{SteelBlue}{HTML}{4682B4}
\definecolor{TealBlue}{HTML}{008080}
\definecolor{DarkSlate}{HTML}{2F4F4F}
\definecolor{DustyGrey}{HTML}{696969}

\makeatletter
\AtBeginDocument{\let\@elt\relax}
\makeatother

\headstyles{bringhurst}
\chapterstyle{dash}

\renewcommand{\chapnamefont}{\normalfont\huge\bfseries\color{DeepNavy}}
\renewcommand{\chapnumfont}{\normalfont\huge\bfseries\color{DeepNavy}}
\renewcommand{\chaptitlefont}{\normalfont\Huge\bfseries\color{DeepNavy}}

\counterwithin{footnote}{chapter}

\setsubparaheadstyle{\normalfont\normalsize\itshape\color{DustyGrey}}
\setlength{\beforechapskip}{-20pt} 
\setlength{\afterchapskip}{30pt}
\setsecnumformat{\csname the#1\endcsname.\quad}

%-----------------------------------------------------------------------------
% Standalone Chapter vs Master Book Modes
%-----------------------------------------------------------------------------
\newcommand{\Copyright}{This work is licensed under the Creative Commons Attribution-Noncommercial 3.0 License. To view a copy of this license, visit \url{http://creativecommons.org/licenses/by-nc/3.0/} or send a letter to Creative Commons, 171 Second Street, Suite 300, San Francisco, California, 94105, USA.}

\newcommand{\RealChapterEndInner}{%
   {%
      \begingroup%
      \emergencystretch=1em%
      \printbibliography[heading=subbibliography,segment=\therefsegment]{}%
      \endgroup%
   }%
}

\providecommand{\BookMode}[1]{}
\providecommand{\StandAloneMode}[1]{#1}

\ifx\bookMode\undefined
   % STANDALONE CHAPTER MODE
   \renewcommand{\StandAloneMode}[1]{#1}
   \renewcommand{\BookMode}[1]{}

   \newcommand{\ChapterNumPage}[2]{%
      \if\relax\detokenize{#1}\relax
         \setcounter{chapter}{0}%
      \else\ifnum#1>0
         \setcounter{chapter}{\numexpr#1-1\relax}%
      \else
         \setcounter{chapter}{0}%
      \fi\fi
      \if\relax\detokenize{#2}\relax
         \setcounter{page}{1}%
      \else\ifnum#2>0
         \setcounter{page}{#2}%
      \else
         \setcounter{page}{1}%
      \fi\fi
   }

   \DeclareRobustCommand{\setheadingcolor}{\color{RoyalBlue}}
   \setsecheadstyle{\setheadingcolor\normalfont\Large\bfseries}
   \setsubsecheadstyle{\normalfont\large\bfseries\color{SteelBlue}}
   \setsubsubsecheadstyle{\normalfont\normalsize\bfseries\color{TealBlue}}
   \setparaheadstyle{\normalfont\normalsize\bfseries\color{DarkSlate}}

   \DeclareDocumentCommand{\Chapter}{s o m}{%
      \IfBooleanTF{#1}{%
         \IfValueTF{#2}{\chapter*[]{#2}}{\chapter*[]{#3}}%
      }{%
         \IfValueTF{#2}{\chapter*[]{#2}}{\chapter[]{#3}}%
      }%
      \setcounter{footnote}{0}%
   }
   \newcommand{\ChapterPrefix}{}
   \newcommand{\ChapterGen}{\pagestyle{plain}\vspace*{-0.7cm}}
   \newcommand{\ChapterEnd}{%
      \RealChapterEndInner{}%
      \end{document}%
   }

   \renewcommand{\File}[1]{#1}
   \renewcommand{\RootFile}[1]{../#1}
\else
   % MASTER BOOK MODE
   \renewcommand{\StandAloneMode}[1]{}
   \renewcommand{\BookMode}[1]{#1}

   \newcommand{\ChapterNumPage}[2]{}

   \DeclareDocumentCommand{\Chapter}{s o m}{%
      \IfBooleanTF{#1}{%
         \IfValueTF{#2}{\chapter*[#2]{#3}}{\chapter*{#3}}%
         \ifx\Dir\undefined\else
            \typeout{CHAPTER_MAP: \Dir/\CurrentChapFile => chap: 0, page: \thepage}%
         \fi
      }{%
         \IfValueTF{#2}{\chapter[#2]{#3}}{\chapter{#3}}%
         \ifx\Dir\undefined\else
            \typeout{CHAPTER_MAP: \Dir/\CurrentChapFile => chap: \thechapter, page: \thepage}%
         \fi
      }%
   }
   \newcommand{\ChapterPrefix}{%
      \vspace{-1cm}%
      \subsection*{\small
         {598 - Class notes for Randomized Algorithms}\\
         {Sariel Har-Peled}\\
         {\today}}%
   }
   \newcommand{\ChapterGen}{}
   \newcommand{\ChapterEnd}{\RealChapterEndInner{}}

   \renewcommand{\File}[1]{\Dir/#1}
   \renewcommand{\RootFile}[1]{#1}

   \newcommand{\InputChap}[2]{%
      \def\Dir{#1}%
      \def\CurrentChapFile{#2}%
      \typeout{}%
      \typeout{------------------------------------------------------------}%
      \typeout{PROCESSING FILE: #1/#2}%
      \typeout{}%
      \input{#1/#2}%
   }
\fi

%-----------------------------------------------------------------------------
% Fragments & Quotes
%-----------------------------------------------------------------------------
\newcommand{\InputExt}[1]{%
   \IfFileExists{#1}%
   {\input{#1}}%
   {\IfFileExists{../#1}%
      {\input{../#1}}%
      {\IfFileExists{../../#1}%
         {\input{../../#1}}%
         {\typeout{ERROR: Unable to find #1}}%
      }%
   }%
}
\newcommand{\FragmentName}[1]{\RootFile{fragment/#1.tex}}

\newwrite\FragOutHandle
\newwrite\FragProbeOut
\newread\FragProbeIn
\newif\ifInJunkDir
\InJunkDirfalse

% Detect whether XeTeX was invoked with -output-directory=junk
\edef\FragProbeSecret{\number\uniformdeviate 1000000}
\immediate\openout\FragProbeOut=frag_probe.tmp
\immediate\write\FragProbeOut{\FragProbeSecret}
\immediate\closeout\FragProbeOut

\openin\FragProbeIn=junk/frag_probe.tmp
\ifeof\FragProbeIn
  \InJunkDirfalse
\else
  \read\FragProbeIn to \FragReadSecret
  \closein\FragProbeIn
  \def\trimspace#1 {#1}
  \edef\FragReadSecretClean{\expandafter\trimspace\FragReadSecret}
  \ifx\FragReadSecretClean\FragProbeSecret
    \InJunkDirtrue
  \else
    \InJunkDirfalse
  \fi
\fi

\newcommand{\FragWritePrefix}{%
  \ifx\bookMode\undefined
    % Standalone chapter mode
    \ifInJunkDir ../../fragment/\else ../fragment/\fi
  \else
    % Master book mode
    \ifInJunkDir ../fragment/\else fragment/\fi
  \fi
}

\DeclareDocumentEnvironment{fragment}{m +b}{%
   \immediate\openout\FragOutHandle=\FragWritePrefix#1.tex
   \immediate\write\FragOutHandle{\unexpanded{#2}}%
   \immediate\closeout\FragOutHandle
   #2%
}{}

\makeatletter
\DeclareDocumentCommand{\IncFragment}{m}{%
   \begingroup
      \let\label\@gobble
      \InputExt{\FragmentName{#1}}%
   \endgroup
}
\makeatother

\newcommand{\QuoteOpen}{}
\newcommand{\QuoteClose}{}
\newcommand{\QuoteWrite}[2]{}

\newcommand{\Quote}[2]{%
   \SaveIndent
   \QuoteWrite{#1}{#2}%
   \begin{center}%
      {\footnotesize
         \begin{quote}
             \RestoreIndent
             {#1}\\
             \hspace*{2cm}-- {#2.}
         \end{quote}%
      }%
      \index{quotation!#2}%
   \end{center}%
   \bigskip
}

\newcommand{\QuotePExt}[3]{%
   \hfill%
   \begin{minipage}{0.95\linewidth}%
       \footnotesize {#1}%
       \rule{\linewidth}{0.3mm}%
       \hfill #2, #3%
   \end{minipage}%
}

\epigraphfontsize{\small\itshape}
\setlength{\epigraphwidth}{0.6\textwidth}
\newcommand{\QuoteExt}[3]{%
   \vspace*{-1.6cm}%
   \epigraph{{\footnotesize #1}}{{#2}, {#3}}%
   \vspace*{-1.2cm}%
}

\providecommand{\delX}[1]{{\color{red}\sout{#1}}}
\providecommand{\newX}[1]{{\color{blue}#1}}
\providecommand{\chgY}[2]{{\color{red}\sout{#1}}{\color{blue}#2}}
\providecommand{\remX}[1]{{\color{purple}\textbf{[Remark: #1]}}}
\newcommand{\SpellIgnore}[1]{#1}
\newcommand{\EQ}{\mathrel{\;\;=\;\;}}

%-----------------------------------------------------------------------------
% Save a symbol that we know is going to get redefined.
\def\savesymbol#1{%
  \expandafter\let\expandafter\origsym\expandafter=\csname#1\endcsname
  \expandafter\let\csname orig#1\endcsname=\origsym
  \expandafter\let\csname#1\endcsname=\relax
}

% Restore a previously saved symbol, and rename the current one.
\def\restoresymbol#1#2{%
  \expandafter\let\expandafter\newsym\expandafter=\csname#2\endcsname
  \expandafter\global\expandafter\let\csname#1#2\endcsname=\newsym
  \expandafter\let\expandafter\origsym\expandafter=\csname orig#2\endcsname
  \expandafter\global\expandafter\let\csname#2\endcsname=\origsym
}

\usepackage{pictex}%
\savesymbol{triangle}%
\usepackage{tree}%
\restoresymbol{TreePKG}{triangle}%

% Definitions absorbed from sbook.sty
\providecommand{\MdPrc}{\nu}

\providecommand{\TwoFiguresBase}[4]{%
   \begin{tabular}{cc}
      \begin{minipage}{0.48\linewidth}
          \begin{center}
              {{{#1}}}
          \end{center}
      \end{minipage}
      &
      \begin{minipage}{0.48\linewidth}
          \begin{center}
              {{#3}}
          \end{center}
      \end{minipage}
      \\
      {#2} & {#4}
   \end{tabular}%
}

\providecommand{\TwoFigures}[6]{%
   \begin{tabular}{cc}
      \begin{minipage}{0.48\linewidth}
          \begin{center}
              {\IncludeGraphics[#1]{{#2}}}
          \end{center}
      \end{minipage}
      &
      \begin{minipage}{0.48\linewidth}
          \centerline{\IncludeGraphics[#4]{{#5}}}
      \end{minipage}
      \\
      {#3} & {#6}
   \end{tabular}%
}

\providecommand{\ThreeFigures}[6]{%
   \begin{tabular}{ccc}
        \begin{minipage}{0.30\linewidth}
            \IncludeGraphics[width=5.1cm]{{#1}}\\
        \end{minipage}
        &
        \begin{minipage}{0.30\linewidth}
            \IncludeGraphics[width=5.1cm]{{#2}}\\
        \end{minipage}
        &
        \begin{minipage}{0.30\linewidth}
            \IncludeGraphics[width=5.1cm]{{#3}}\\
        \end{minipage}
        \vspace{-1cm}
        \\
        {#4} & {#5}& {#6}
   \end{tabular}%
}

%-----------------------------------------------------------------------------
\providecommand{\remove}[1]{}
\providecommand{\NotAMSMode}[1]{#1}
\providecommand{\Solution}[1]{}
\providecommand{\Version}[1]{\hfill{Version: #1}}
\providecommand{\Distor}{\mathrm{dist}}
\providecommand{\Hardness}[1]{(h:#1)}
\providecommand{\QuoteOptExt}[3]{\QuoteExt{#1}{#2}{#3}}
\providecommand{\Hypercube}{\EuScript{H}}
\providecommand{\emphiC}[1]{\emphic{#1}{#1}}
\providecommand{\CoverTime}{\EuScript{C}}
\providecommand{\Nbr}{\Gamma}
\providecommand{\CommuteTime}[2]{\mathbf{CT}_{#1#2}}
\providecommand{\Resistance}{\mathbf{R}}
\providecommand{\EffectResist}[2]{\mathbf{R}_{#1#2}}
\providecommand{\NetworkE}{\EuScript{N}}
\providecommand{\trinom}[4]{\binom{#1}{#2 \;\; #3 \;\; #4}}
\providecommand{\Ranges}{\EuScript{R}}
\providecommand{\CClass}[1]{\ensuremath{\mathbf{#1}}}
\providecommand{\coCClass}[1]{\ensuremath{\mathrm{co-}}\CClass{#1}}
\providecommand{\CX}{\EuScript{C}}
\providecommand{\packet}{\mathbf{c}}
\providecommand{\packetA}{\mathbf{d}}
\providecommand{\mopt}{m_{\mathrm{opt}}}
\providecommand{\girth}{\mathop{\mathrm{girth}}}
\providecommand{\CR}{\mathop{\mathrm{cr}}}
\providecommand{\TRUE}{{\textcolor[named]{RedViolet}{\textsf{TRUE}}}\xspace}
\providecommand{\FALSE}{{\textcolor[named]{RedViolet}{\textsf{FALSE}}}\xspace}
\providecommand{\Part}[1]{\part{#1}}
\providecommand{\Section}[1]{\section{#1}}
\providecommand{\SubSection}[1]{\subsection{#1}}
\providecommand{\SubSubSection}[1]{\subsubsection{#1}}
\providecommand{\Exercise}[5]{%
\begin{exercise_h}[#1]%
   \exelab{#3}{}%
   #4%
\end{exercise_h}%
}
\providecommand{\si}[1]{#1}
\providecommand{\myParagraph}[1]{\noindent\textbf{#1}}
\providecommand{\DefName}[1]{#1}
\providecommand{\deflabel}[1]{def_#1}
\providecommand{\Description}[1]{}

%-----------------------------------------------------------------------------
% Mathematical Notations & Environments (from notations.sty)
%-----------------------------------------------------------------------------
%-----------------------------------------------------------------------------
% List environments (questions & compactenum)
%-----------------------------------------------------------------------------

\newlist{compactenumA}{enumerate}{5}%
\setlist[compactenumA]{topsep=0pt,itemsep=-1ex,partopsep=1ex,parsep=1ex,%
   label=(\Alph*)}%

\newlist{compactenumI}{enumerate}{5}%
\setlist[compactenumI]{topsep=0pt,itemsep=-1ex,partopsep=1ex,parsep=1ex,%
   label=(\Roman*)}%

\newlist{compactenuma}{enumerate}{5}%
\setlist[compactenuma]{topsep=0pt,itemsep=-1ex,partopsep=1ex,parsep=1ex,%
   label=(\alph*)}%

\newlist{compactitem}{enumerate}{5}%
\setlist[compactitem]{topsep=0pt,itemsep=-1ex,partopsep=1ex,parsep=1ex,%
   label=\textbullet}%

\newlist{compactenumi}{enumerate}{5}%
\setlist[compactenumi]{topsep=0pt,itemsep=-1ex,partopsep=1ex,parsep=1ex,%
   label=(\roman*)}%

\setcounter{secnumdepth}{5}

\newcommand{\HST}{\TermI{HST}\xspace}
\newcommand{\Term}[1]{\textsf{#1}}
\newcommand{\TermI}[1]{\Term{#1}\index{#1@\Term{#1}}}

\DeclareMathAlphabet{\mathpzc}{OT1}{pzc}{m}{it}

\newcommand{\seg}{\mathsf{s}}

\newcommand{\Diam}{\mathrm{diam}}
\newcommand{\diam}{\mathrm{diam}}

\newcommand{\Copt}{C_{\mathrm{opt}}}
\newcommand{\CenterPnt}{\mathrm{cen}}

\newcommand{\parent}{\overline{\mathrm{p}}}
\newcommand{\parentX}[2][\!]{\parent\pth[#1]{#2}}

\newcommand{\FMS}{\Mh{\EuScript{X}}\index{metric space}}  
\newcommand{\FMSB}{\Mh{\EuScript{Y}}\index{metric space}}  
\newcommand{\FastCut}{\Algorithm{FastCut}\xspace}
\newcommand{\MinCut}{\Algorithm{MinCut}\xspace}
\newcommand{\MinCutRep}{\Algorithm{MinCutRep}\xspace}
\newcommand{\Contract}{\Algorithm{Contract}\xspace}

\newcommand{\Grid}{\mathrm{G}}
\newcommand{\ClosestPair}{\mathcal{CP}}
\newcommand{\SphereChar}{\ensuremath{\mathbb{S}}}
\newcommand{\Sphere}[1]{{\SphereChar^{(#1)}}}

\newcommand{\PntSet}{{\Mh{\mathsf{P}}}}%
\providecommand{\PntSetA}{{\Mh{\mathsf{Q}}}}%

\newcommand{\query}{\mathcalb{q}}

\newcommand{\pnt}{\mathsf{p}}
\newcommand{\pntA}{\mathsf{u}}
\newcommand{\pntB}{\mathsf{v}}
\newcommand{\pntC}{{\mathsf{x}}}
\newcommand{\pntD}{{\mathsf{y}}}
\newcommand{\pntE}{{\mathsf{z}}}

\newcommand{\pntCoord}{\Mh{p}}
\newcommand{\pntCoordA}{\Mh{q}}

%%%%%%%%%%%%%%%%%%%%%%%%%%%%%%%%%%%%%%%%%%%%%%%%%%%%%%%%%%%%%%%%%%%%% 
%%%%%%%%%%%%%%%%%%%%%%%%%%%%%%%%%%%%%%%%%%%%%%%%%%%%%%%%%%%%%%%%%%%%%
% Markov chain notations.

\newcommand{\MCSpc}{\Mh{\mathtt{S}}}
\newcommand{\MCPrb}{\Mh{\mathtt{P}}}
\newcommand{\MCPrbY}[2]{\MCPrb_{#1#2}}
\newcommand{\MCPrbM}{\Mh{\mathbf{P}}}%
\newcommand{\MCPrbT}[3]{\MCPrb_{#1#2}^{(#3)}}
\newcommand{\MCFPrbT}[3]{\Mh{\mathtt{r}}_{#1#2}^{(#3)}}
\newcommand{\MCVstProb}[2]{\Mh{\mathtt{f}}_{#1#2}}
\newcommand{\MCExSteps}[2]{\Mh{\mathtt{h}}_{#1#2}}
\newcommand{\MCPVec}[1]{\mathbf{q}^{(#1)}}
%%%%%%%%%%%%%%%%%%%%%%%%%%%%%%%%%%%%%%%%%%%%%%%%%%%%%%%%%%%%%%%%%%%%%

\newcommand{\Spread}{\Phi}
\newcommand{\VC}{\ensuremath{\mathsf{VC}}\xspace}
\newcommand{\MC}{\ensuremath{\mathsf{MC}}\xspace}
\newcommand{\hyperplane}{h}

\newcommand{\QuickSort}{\Algorithm{QuickSort}\xspace}
\newcommand{\Alg}{\Algorithm{Alg}\xspace}

\newcommand{\Graph}{\Mh{\mathsf{G}}}%
\newcommand{\GraphB}{\mathsf{H}}
\newcommand{\GraphC}{\mathsf{K}}
\newcommand{\GraphD}{\mathsf{F}}

\providecommand{\G}{}
\renewcommand{\G}{\Graph}

\newcommand{\RLP}{\ComplexityClass{RLP}\xspace}
\newcommand{\USTCON}{\ComplexityClass{USTCON}}
\newcommand{\STCON}{\overrightarrow{\ComplexityClass{STCON}}}

\newcommand{\Mtrx}{\mathsf{M}}
\newcommand{\TMtrx}{\mathsf{Q}}
\newcommand{\IMtrx}{\mathcal{I}}

\newcommand{\Expander}[3]{$\left[#1,#2,#3\right]$-expander}

\newcommand{\ChartX}[1]{\ensuremath{\chi_{#1}}}
\newcommand{\Field}{\mathbb{F}}
\newcommand{\LSpace}{\mathbb{G}}
\newcommand{\LD}[2]{\mathsf{L{D}}\pth{#1,#2}}
\newcommand{\SEVN}{\widehat{\lambda_2}}
\newcommand{\NX}[2]{#1\!\left[#2\right]}
\newcommand{\Entry}[2]{#1\!\left[#2\right]}
\newcommand{\NXExt}[3]{#1_{#2}\!\left[#3\right]}
\newcommand{\RangeI}[1]{\llbracket #1 \rrbracket}
\newcommand{\DialX}[1]{\gamma\pth{#1}}
\newcommand{\BiDialX}[1]{{\gamma}_2\pth{#1}}

\newcommand{\EV}[1]{\widehat{\lambda_{#1}}}
\newcommand{\TEV}[1]{\widehat{\lambda_{#1}}}
\newcommand{\EVal}{\widehat{\lambda}}
\newcommand{\Basis}{\mathcal{B}}

\newcommand{\ZProd}[2]{\ensuremath{#1 \, \text{\textcircled{z}} #2}}
\newcommand{\RProd}[2]{\ensuremath{#1 \, \text{\textcircled{r}} #2}}

\newcommand{\OneVec}[1]{\mathbf{1}^{#1}}
\newcommand{\DstrbV}[1]{\mathbf{\mathsf{q}}^{(#1)}}
\newcommand{\DstrbE}[2]{\mathsf{q}^{(#1)}_{#2}}

\newcommand{\DEdges}{\widetilde{\Edges}}

\newcommand{\VV}{\Mh{\mathsf{V}}}%
\newcommand{\EE}{\Mh{\mathsf{E}}}%

\newcommand{\Vertices}{\VV}%
\newcommand{\Edges}{\EE}%

\newcommand{\EdgesX}[1]{\Edges\pth{#1}}
\newcommand{\VerticesX}[1]{V\pth{#1}}

\newcommand{\StringX}[1]{\mathsf{s}\pth{#1}}

\newcommand{\BMtrx}{\mathsf{B}}
\newcommand{\WMtrx}{\mathsf{W}}
\newcommand{\NWMtrx}{\overline{\mathsf{W}}}
\newcommand{\Sequence}{\EuScript{S}}

\newcommand{\med}{\mathrm{med}}
\newcommand{\medX}[1]{\Mh{\mathrm{med}}\pth{#1}}

\newcommand{\medianC}{\Mh{\nu}}%

\newlength{\savedparindent}
\newcommand{\SaveIndent}{\setlength{\savedparindent}{\parindent}}
\newcommand{\RestoreIndent}{\setlength{\parindent}{\savedparindent}}

\newcommand{\NNN}{\mathbb{N}}%
\providecommand{\NN}{\mathbb{N}}%
\renewcommand{\Re}{\mathbb{R}}%
\newcommand{\reals}{\Re}%
\newcommand{\ZZ}{\mathbb{Z}}%

\newcommand{\level}{\mathop{\mathrm{level}}}

\providecommand{\TPDF}[2]{\texorpdfstring{#1}{#2}}

\newcommand{\Hint}[1]{\emph{[#1]}}
\newcommand{\Arr}{\mathop{\mathrm{\EuScript{A}}}}

\newcommand{\Array}{\mathcal{T}}%

\newcommand{\dirEdge}[2]{\left({{#1} \rightarrow {#2}} \right)}

%%%%%%%%%%%%%%%%%%%%%%%%%%%%%%%%%%%%%%%%%%%%%%%%%%%%%%%%%%%%%%%%%%%%%%
% Emphasize stuff...
\definecolor{blue25}{rgb}{0, 0, 11}
\newcommand{\emphic}[2]{%
   \textcolor{blue25}{%
      \bfseries{\emph{#1}}}%
   \index{#2}}
\newcommand{\emphi}[1]{\emphic{#1}{#1}}
\newcommand{\emphOnly}[1]{\textcolor{blue25}{\bfseries{\emph{#1}}}}

\definecolor{almostblack}{rgb}{0, 0, 0.3}
\providecommand{\emphw}[1]{{\textcolor{almostblack}{\emph{#1}}}}%

%
%%%%%%%%%%%%%%%%%%%%%%%%%%%%%%%%%%%%%%%%%%%%%%%%%%%%%%%%%%%%%%%%%%%%%%

\newcommand{\volX}[1]{\mathrm{vol}\pth{#1}}

\newcommand{\Interval}{\EuScript{I}}
\newcommand{\IntervalA}{\EuScript{I}}

\newcommand{\Exp}{\mathrm{Exp}}

\let\Ball\undefined
\newcommand{\Ball}{\mathbf{b}}
\newcommand{\hdim}{d}
\newcommand{\Domain}{\mathcal{D}}
\newcommand{\points}[1]{[\bfseries{#1 points}]}
\newcommand{\Event}{\Mh{\EuScript{E}}}%
\newcommand{\EventB}{\Mh{\EuScript{B}}}%
\newcommand{\EventF}{\Mh{\EuScript{F}}}%
\newcommand{\EventG}{\Mh{\EuScript{G}}}%

\newcommand{\distCmd}[1]{\left\| {#1}  \right\|}
\newcommand{\distY}[2]{\distCmd{#1 - #2}}

\newcommand{\NormalC}{\mathbf{N}}
\newcommand{\NormalY}[2]{\NormalC\pth{#1, #2}}%
\newcommand{\NormalDistOneDim}{\NormalC}
\newcommand{\NormalDistX}[1]{\NormalC^{#1}}
\DeclareMathOperator{\lca}{lca}
\newcommand{\RangeSpace}{\Mh{\mathsf{S}}}%
\newcommand{\RangeSpaceA}{\mathsf{T}}

\newcommand{\range}{\mathbf{r}}
\newcommand{\rangeA}{\mathbf{r}'}

\newcommand{\BadProb}{\varphi}

\newcommand{\EntropyChar}{\mathbb{H}}
\newcommand{\Entropy}[1]{\EntropyChar\pth{#1}}
\newcommand{\EntropyX}[1]{\EntropyChar\pth[]{#1}}

\newcommand{\Matrix}{\mathsf{M}}
\newcommand{\MatrixB}{\mathsf{B}}
\newcommand{\MatrixC}{\mathsf{C}}
\newcommand{\MatrixD}{\mathsf{D}}

\renewcommand{\th}{th\xspace}

\newcommand{\Family}{\mathcal{F}}
\newcommand{\FamilyB}{\mathcal{G}}
\newcommand{\FamilyC}{\mathcal{H}}
\newcommand{\FMSC}{\EuScript{Y}\index{metric space}}

\newcommand{\median}{\overline{\mathsf{m}}}
\newcommand{\QuickSelect}{\Algorithm{Quick{}Select}\xspace}

\providecommand{\row}{\mathbf{r}}%
\newcommand{\rowc}{{r}}%

\newcommand{\vecb}{\mathbf{b}}

\newcommand{\vecA}{\mathbf{v}}
\newcommand{\vecB}{\mathbf{u}}
\newcommand{\vecRand}{\mathbf{r}}

\newcommand{\vecH}{\mathbf{h}}
\newcommand{\vecV}{\mathbf{v}}

\newcommand{\vecX}{\mathbf{x}}
\newcommand{\vecY}{\mathbf{y}}

\newcommand{\DistBinom}[2]{\mathrm{Bin}\pth{#1,#2}}

\newcommand{\DistGeom}[1]{\mathrm{Geom}\pth{#1}}

\newcommand{\Sample}{\mathsf{R}}

\newcommand{\lcm}{\mathop{\mathrm{lcm}}}
\newcommand{\bdiv}{\mathop{\mathrm{div}}}%
\newcommand{\divides}{\mid}%
\newcommand{\poly}{\mathrm{poly}}%

\newcommand{\Group}{\mathcal{G}} 
\newcommand{\GroupA}{\mathcal{H}} 

\newcommand{\algGCD}{\Algorithm{EuclidGCD}\xspace}%

\newcommand{\gOp}{\times}%
\newcommand{\gIdent}{\mathsf{i}}%

\newcommand{\ordX}[2][\!]{\mathop{\mathrm{o{r}d}}\pth[#1]{#2}}

\newcommand{\lSym}[2]{%
   \left(
       \ds {#1}\divides{#2}%
   \right) %
}
\newcommand{\jSym}[2]{\left\llbracket%
       \ds {#1}\divides{#2}%
       \right\rrbracket%
}
\newcommand{\Line}{\ell}

\newcommand{\LevelChar}{\mathbb{L}}
\newcommand{\AtMostKLevelChar}{\LevelChar_{\leq k}}
\newcommand{\AtMostKLevel}[1]{\LevelChar_{\leq k} \pth{ #1}}

\newcommand{\LineSet}{\mathsf{L}}%
\newcommand{\ArrX}[1]{\EuScript{A}\pth{#1}}%

\providecommand{\Matousek}{Matou{\v s}ek\xspace}
\providecommand{\Erdos}{Erd{\H o}s\xspace}
\providecommand{\Lovasz}{Lov\'asz\xspace}

\newcommand{\ProblemC}[1]{{\textsf{\textcolor{red}{#1}\index{problem!#1}}}}
\DefineNamedColor{named}{OliveGreen} {cmyk}{0.64,0,0.95,0.40}
\providecommand{\ComplexityClass}[1]{{{\textcolor[named]{OliveGreen}{%
            \textsc{#1}}}}}
\providecommand{\NPHard}{{\ComplexityClass{NP-Hard}}%
   \index{NP!hard}\xspace}
\providecommand{\NPComplete}{\ComplexityClass{NP-Complete}%
   \index{NP!complete}\xspace}
\providecommand{\APXHard}{\ComplexityClass{APX-Hard}%
   \index{APX!Hard}\xspace}
\providecommand{\NPC}{\ComplexityClass{NPC}%
   \index{NP!complete}\xspace}
\providecommand{\NP}{\ComplexityClass{NP}%
   \index{NP}\xspace}
\providecommand{\POLYT}{\ComplexityClass{P}\xspace}

\newcommand{\envFrame}[1]{\begin{center}\fbox{\begin{minipage}[c]{0.75\textwidth}{#1}\end{minipage}}\end{center}}

\newcommand{\mysqrt}[1]{\ts\!\!\sqrt{#1}}
\newcommand{\SegSet}{\mathsf{S}}
\newcommand{\ArrVD}[1]{\Arr^{|} \pth{#1}}
\newcommand{\RIC}{\textsf{RIC}\index{random incremental
      construction}\xspace}
\newcommand{\ArrB}{\EuScript{B}}

\newcommand{\TreeOpt}{\EuScript{T}_\mathrm{opt}}
\newcommand{\TreeOptA}{\EuScript{T}_\mathrm{opt}'}

\newcommand{\DepthExt}[2]{\mathrm{depth}_{#1}\pth{#2}}
\newcommand{\depthX}[1]{\mathrm{depth}\pth{#1}}

\newcommand{\TreeA}{\EuScript{U}}
\newcommand{\TreeB}{\EuScript{V}}

\newcommand{\Trap}{\sigma}
\newcommand{\TrapA}{\sigma'}
\newcommand{\TrapB}{\sigma''}

\newcommand{\SampleA}{\mathsf{R}'}

\newcommand{\AllRegions}{\mathcal{T}}
\newcommand{\RegionsSet}[1]{\mathcal{F}\pth{#1}}
\newcommand{\RegionsHeavy}[2]{\mathcal{F}_{\!\geq#1}\pth{#2}}
\newcommand{\NumOfRegionsZero}[1]{\ExChar \!f\pth{#1}}
\newcommand{\NumOfRegions}[2]{\ExChar\! f_{\!\geq #1}\pth[]{#2}}
\newcommand{\ObjSet}{\mathsf{S}}

\newcommand{\WeightX}[2][\!]{\omega \pth[#1]{#2}}
\newcommand{\Edge}{\mathsf{e}}

\newcommand{\Face}{\mathsf{f}}

\newcommand{\DefSet}[1]{D\pth{#1}}
\newcommand{\KillSet}[1]{K\pth{#1}}

\newcommand{\CSProbFull}[4]{
   \,\rho_{\!\!\stackrel{\rule[-.005cm]{0cm}{0.1cm}}{#1,
       #4}}\pth{\MakeSBig #2,#3}}

\newcommand{\disk}{D}
\newcommand{\diskA}{D'}

\newcommand{\IntervalSet}{\widehat{\mathcal{I}}}

\newcommand{\GroundSet}{\textsf{X}}
\newcommand{\FGroundSet}{\mathsf{x}}
\newcommand{\RangeSet}{{\mathcal{R}}}
\newcommand{\RangeSetA}{\mathcal{R}'}
\newcommand{\MeasureChar}{\overline{m}}
\newcommand{\Measure}[1]{\MeasureChar\pth{#1}}
\newcommand{\sMeasure}[1]{\overline{s}\pth{#1}}

\newcommand{\SetA}{\Mh{S}}%
\newcommand{\SetB}{\Mh{T}}%
\newcommand{\SetC}{\Mh{N}}%

\newcommand{\VCProj}[2]{#1_{|#2}}
\newcommand{\VCDim}{\mathrm{dim}_{\VC}}
\newcommand{\Dim}{\delta}
\newcommand{\DimA}{\delta'}
\newcommand{\PntSetB}{{{\mathsf{Q}}}}
\newcommand{\PntSetC}{{{\mathsf{R}}}}
\newcommand{\Growth}[2]{{\EuScript{G}}_{#1}\pth{#2}}
\newcommand{\ShatterFunc}[1]{\pi_{#1}}
\newcommand{\Distribution}{\EuScript{D}}
\newcommand{\Error}{\EuScript{E}}

\newcommand{\bd}{{\partial}}%
\newcommand{\eps}{{\varepsilon}}%

\newcommand{\clmlab}[1]{\label{claim:#1}}
\newcommand{\clmref}[1]{\HLink{Claim}{claim:#1}}

\newcommand{\CHX}[1]{{\mathcal{CH}}\pth{#1}}

\newcommand{\ts}{\hspace{0.6pt}}

\def\qedsymbol{\myqedsymbol}%

\newcommand{\qedsign}{\hfill{\hfill\myqedsymbol}}
\newcommand{\qedsignBig}{\hfill{\hfill\rule{4mm}{4mm}}}
\newcommand{\QED}{{\hfill{\hfill\qedsymbol}}}

\newcommand{\num}{\alpha}

\newcommand{\NEW}[1]{#1}

\newcommand{\ThreeFiguresExt}[9]{%
   \begin{tabular}{ccc}
     \begin{minipage}{0.31\linewidth}
         \begin{center}
             {\IncludeGraphics[#1]{{#2}}}
         \end{center}
     \end{minipage}
     &
       \begin{minipage}{0.31\linewidth}
           \centerline{\IncludeGraphics[#4]{{#5}}}
       \end{minipage}
     &
       \begin{minipage}{0.31\linewidth}
           \centerline{\IncludeGraphics[#7]{{#8}}}
       \end{minipage}
     \\
     ~\\[-0.2cm]
     {#3} & {#6} & {#9}
   \end{tabular}}

\newcommand{\IncludeGraphics}[2][]{%
   \typeout{}%
   \typeout{Graphics: #2}%
   \typeout{\ includegraphics[#1]{#2}}%
   \includegraphics[#1]{#2}
   \typeout{}%
}

%    -----------------------------------------------------------------
%    -----------------------------------------------------------------
%    -----------------------------------------------------------------
%    Handling references
%    -----------------------------------------------------------------
%    -----------------------------------------------------------------
%    -----------------------------------------------------------------
\newcommand{\HLinkShort}[2]{\hyperref[#2]{#1\ref*{#2}}}
\newcommand{\HLink}[2]{\hyperref[#2]{#1~\ref*{#2}}}
\newcommand{\HLinkPage}[2]{\hyperref[#2]{#1~\ref*{#2}%
      $_\text{p\pageref{#2}}$}}

\newcommand{\HLinkSuffix}[3]{\hyperref[#2]{#1\ref*{#2}{#3}}}
\newcommand{\HLinkPageSuffix}[3]{\hyperref[#2]{#1\ref*{#2}%
      #3$_\text{p\pageref{#2}}$}}

\newcommand{\figlab}[1]{\label{fig:#1}}
\newcommand{\figref}[1]{\HLink{Figure}{fig:#1}}

\newcommand{\exelab}[1]{\label{exer:#1}}
\newcommand{\exeref}[1]{Exercise~\ref{exer:#1}}

\providecommand{\deflab}[1]{\label{def:#1}}
\newcommand{\defref}[1]{\HLink{Definition}{def:#1}}
\newcommand{\defrefY}[2]{\hyperref[def:#2]{#1}}
\newcommand{\defrefpage}[1]{\HLinkPage{Definition}{def:#1}}

\newcommand{\etal}{\textit{et~al.}\xspace}

\newcommand{\obslab}[1]{\label{observation:#1}}
\newcommand{\obsref}[1]{\HLink{Observation}{observation:#1}}
\newcommand{\obsrefY}[2]{\hyperref[observation:#2]{#1}}

\newcommand{\exmlab}[1]{\label{example:#1}}
\newcommand{\exmref}[1]{\HLink{Example}{example:#1}}

\newcommand{\remlab}[1]{\label{rem:#1}}
\newcommand{\remref}[1]{\HLink{Remark}{rem:#1}}%

\newcommand{\chaplab}[1]{\label{chap:#1}}

\newcommand{\itemlab}[1]{\label{item:#1}}
\newcommand{\itemref}[1]{\HLinkSuffix{}{item:#1}{}}

\newcommand{\tedlab}[1]{\label{tedium:#1}}
\newcommand{\tedref}[1]{\HLink{Tedium}{tedium:#1}}%

\newcommand{\factlab}[1]{\label{fact:#1}}%

\newcommand{\alglab}[1]{\label{algorithm:#1}}

\newcommand{\net}{\Mh{\mathcal{N}}}%
      
\newcommand{\lemlab}[1]{\label{lemma:#1}}
\newcommand{\lemref}[1]{\HLink{Lemma}{lemma:#1}}%
\newcommand{\lemrefY}[2]{\hyperref[lemma:#2]{#1}}
\newcommand{\lemrefpage}[1]{\HLinkPage{Lemma}{lemma:#1}}
\newcommand{\lemrefshort}[1]{\HLinkShort{L}{lemma:#1}}

\newcommand{\eqlab}[1]{\label{equation:#1}}
\newcommand{\Eqrefbare}[1]{\HLinkSuffix{(}{equation:#1}{)}}
\newcommand{\Eqref}[1]{\HLinkSuffix{Eq.~(}{equation:#1}{)}}
\newcommand{\Eqrefpage}[1]{\HLinkPageSuffix{Eq.~(}{equation:#1}{)}}

\newcommand{\seclab}[1]{\label{sec:#1}}
\newcommand{\secref}[1]{\HLink{Section}{sec:#1}}
\newcommand{\secrefY}[2]{\hyperref[sec:#2]{#1}}
\newcommand{\secrefpage}[1]{\HLinkPage{Section}{sec:#1}}

\newcommand{\thmlab}[1]{{\label{theo:#1}}}
\newcommand{\thmref}[1]{\HLink{Theorem}{theo:#1}}
\newcommand{\thmrefY}[2]{\hyperref[theo:#2]{#1}}
\newcommand{\thmrefpage}[1]{\HLinkPage{Theorem}{theo:#1}}

\newcommand{\corlab}[1]{\label{cor:#1}}
\newcommand{\corref}[1]{\HLink{Corollary}{cor:#1}}%

%    -----------------------------------------------------------------
%    -----------------------------------------------------------------
%    -----------------------------------------------------------------
%    Handling references
%    -----------------------------------------------------------------
%    -----------------------------------------------------------------
%    -----------------------------------------------------------------

\newcommand{\ceil}[1]{\left\lceil {#1} \right\rceil}
\newcommand{\floor}[1]{\left\lfloor {#1} \right\rfloor}
\newcommand{\brc}[1]{\left\{ {#1} \right\}}
\newcommand{\pth}[2][\!]{\mleft({#2}\mright)}%
\newcommand{\brk}[1]{\mleft[{#1}\mright]}%
\newcommand{\pbrc}[2][\!\!]{#1\left[ {#2} \bigr. \right]}
\newcommand{\pbrcS}[1]{\left[ {#1} \right]}

\newcommand{\dotprod}[2]{{#1}\cdot{#2}}
\newcommand{\DotProd}[2]{\permut{{#1},{#2}}}
\newcommand{\permut}[1]{\left\langle {#1} \right\rangle}

\newcommand{\MakeBig}{}
\newcommand{\MakesBig}{}
\newcommand{\MakeSBig}{}
\newcommand{\MakeVBig}{}

\newcommand{\DistChar}{\Mh{\mathsf{d}}}%
\newcommand{\DistSetY}[2]{\DistChar\pth{#1,#2}}
\newcommand{\distG}[2]{\DistChar_\Graph\pth{#1, #2}}
\newcommand{\DistMtrCharX}[1]{\DistChar_{#1}}

\newcommand{\dY}[2]{\DistChar\pth{#1, #2}}%
\newcommand{\dGY}[2]{\DistChar_\Graph\pth{#1, #2}}%

\newcommand{\cset}{\mathcal{C}}%

\newcommand{\MTR}{\mathcal{M}}
\newcommand{\DistMtr}[2]{\DistChar\pth{#1,#2}}

\newcommand{\dist}[1]{\left\| {#1} \right\|}
\newcommand{\norm}[1]{\left\| {#1} \right\|}

\providecommand{\ds}{\displaystyle}
\newcommand{\sep}[1]{\,\left|\, {#1} \bigr.\right.}%
\newcommand{\sepw}[1]{\left|\, {#1} \right.}%

\newcommand{\ProbExt}[2]{\mathop{\mathbf{Pr}}_{#1}\pbrcx{#2}}
\newcommand{\ProbLTR}{\mathbb{P}}%

\newcommand{\Prob}[1]{\mathop{\ProbLTR} \mleft[ #1 \mright]}%
\newcommand{\ProbCond}[2]{\mathop{\ProbLTR}\!\left[%
       #1 \;\middle\vert\; #2 \right]}
\newcommand{\ProbY}[2]{\ProbLTR^{}_{\!#1}\pbrc{#2}}

\newcommand{\ExChar}{\mathbb{E}}%
\newcommand{\ExSym}{\mathop{\ExChar}}%
\newcommand{\Ex}[2][\!]{\ExSym#1\pbrcx{#2}}
\newcommand{\ExExt}[2]{\ExSym_{#1}\!\pbrcx{#2}}
\newcommand{\ExExtS}[2]{\ExSym_{\substack{#1}}\!\pbrcx{#2}}
\newcommand{\ExCond}[2]{\ExSym\!\left[%
       #1 \;\middle\vert\; #2 \right]}

\newcommand{\VarSymb}{\mathop{\mathbb{V}}}
\newcommand{\Var}[1]{\mathop{\mathbb{V}}\!\pbrcx{#1}}

\newcommand{\pbrcx}[1]{\left[ {#1} \right]}%
\newcommand{\nfrac}[2]{{#1}/{#2}}
\newcommand{\cardin}[1]{\left| {#1} \right|}%

\newcommand\CCancel[2][black]{{\renewcommand\CancelColor{\color{#1}}\cancel{#2}}}

\newcommand{\IncGraphPage}[4][]{%
   \IncludeGraphics[page=#4,#1]{\File{#2/#3}}%
}

\newenvironment{enumEx}{
   \begin{compactenumA}
    }{\end{compactenumA}}

\newcommand{\cl}{{\mathrm{cl}}}

\usepackage[amsmath,thmmarks]{ntheorem}
\theoremseparator{.}%

\newtheorem{theorem:unnumbered}{Theorem}[section]
\newtheorem{theorem}{Theorem}[section]

\newtheorem{observation}[theorem]{Observation}
\newtheorem{conjecture}[theorem]{Conjecture}
\newtheorem{lemma}[theorem]{Lemma}
\newtheorem{corollary}[theorem]{Corollary}
\newtheorem{claim}[theorem]{Claim}
\newtheorem{assumption}[theorem]{Assumption}
\newtheorem{tedium}[theorem]{Tedium}

\theoremstyle{plain}%
\theoremheaderfont{\sffamily}
\theorembodyfont{\upshape}%

\newtheorem{defn}[theorem]{Definition}
\newtheorem{definition}[theorem]{Definition}

\newtheorem*{remark_unnumbered}[theorem:unnumbered]{Remark}%
\newtheorem*{remarks}[theorem]{Remarks}%
\newtheorem{remark}[theorem]{Remark}%
\newtheorem{fact}[theorem]{Fact}%
\newtheorem{example}[theorem]{Example}
\newtheorem{exercise}[theorem]{Exercise}
\newtheorem{problem}[theorem]{Problem}
\newtheorem{algorithm}[theorem]{Algorithm}%
\newtheorem{xca}[theorem]{Exercise}
\newtheorem{openproblem}[theorem]{Open Problem}
\newtheorem{exercise_h}[theorem]{Exercise}

% Proof environments
\theoremheaderfont{\em}%
\theorembodyfont{\upshape}%
\theoremstyle{nonumberplain}
\theoremseparator{}
\theoremsymbol{\myqedsymbol}
\newtheorem{proof}{Proof:}
\newtheorem{proofextin}{Proof}
\newtheorem{altproof}{Alternative proof:}

\theoremsymbol{}
\newtheorem{proofnoqed}{Proof:}

\theoremsymbol{\myqedsymbol}
\theorempreskip{0pt}%
\theorempostskip{0pt}%
\newtheorem{proofof}{Proof of\!}%

\newenvironment{proofExt}[1]{{\emph{#1:}}}{\hfill{\hfill\myqedsymbol}}
\newenvironment{proofnoproof}{{}}{{\hfill{\hfill\myqedsymbol}}}

%%%%%%%%%%%%%%%%%%%%%%%%%%%%%%%%%%%%%%%%%%%%%%%%%%%%%%%%%%%%%%%%%%%%
%%%%%%%%%%%%%%%%%%%%%%%%%%%%%%%%%%%%%%%%%%%%%%%%%%%%%%%%%%%%%%%%%%%%
% Macros to embed figures in text

\newlength{\RecordedWidth}
\newcommand{\RecordWidth}[1]{\settowidth{\RecordedWidth}{#1}}%

\newlength{\msize}%

%-------------------------------------------------------------------
% \minipageWFigure{ TEXT }{ Figure }{ Figure caption/label }
%-------------------------------------------------------------------
\newcommand{\minipageWFigure}[3]{%
   \RecordWidth{#2}%
   \setlength{\msize}{\dimexpr(\textwidth-1.1\RecordedWidth)\relax}
   \noindent%
   \SaveIndent%
   \begin{minipage}{\msize}%
       \RestoreIndent%
       #1
   \end{minipage}%
   \null\hfill%
   \begin{minipage}{\RecordedWidth}
       \null\hfill%
       #2%       

       #3%
   \end{minipage}%
}

\newcommand{\EBRY}[2]{#1_{\langle#2\rangle}}%

\newcommand{\NiceFrac}[2]{%
   {\raisebox{0.2em}{\ensuremath{#1}}\!\!\!}%
   {\rotatebox[origin=c]{-20}{\Big/}}% Rotated /
   {\! \raisebox{-0.2em}{\ensuremath{#2}} }}%

% end - Macros to embed figures in text
%%%%%%%%%%%%%%%%%%%%%%%%%%%%%%%%%%%%%%%%%%%%%%%%%%%%%%%%%%%%%%%%%%%% 
%%%%%%%%%%%%%%%%%%%%%%%%%%%%%%%%%%%%%%%%%%%%%%%%%%%%%%%%%%%%%%%%%%%%

%%%%%%%%%%%%%%%%%%%%%%%%%%%%%%%%%%%%%%%%%%%%%%%%%%%%%%%%%%%%%%%%%%
% Color in math
\providecommand{\Mh}[1]{{#1}}%
\providecommand{\IfColorMath}[2]{#2}%

\IfFileExists{.latex_printer_friendly}{\def\GenPrinterVer{1}}{}%

\ifx\GenPrinterVer\undefined
   \IfFileExists{.latex_color}{\def\GenColorMath{1}}{}
\else
   \renewcommand{\IfPrinterVer}[2]{#1}%
\fi

\ifx\GenColorMath\undefined
\else
\renewcommand{\IfColorMath}[2]{#1}%
\renewcommand{\Mh}[1]{{\textcolor{red}{#1}}}%
\fi
% Color in math - end
%%%%%%%%%%%%%%%%%%%%%%%%%%%%%%%%%%%%%%%%%%%%%%%%%%%%%%%%%%%%%%%%%%

\newcommand{\HSetY}[2]{\Mh{H}\pth{#1, #2}}%
\newcommand{\CHSetY}[2]{\Mh{C}\pth{#1, #2}}%
\newcommand{\Holder}{H\"older\xspace}
\newcommand{\OldOm}{{\mathcal{O}}}
\newcommand{\NewOm}{{\mathcal{N}}}

%%%%%%%%%%%%%%%%%%%%%%%%%%%%%%%%%%%%%%%%%%%%%%%%%%%%%%%%%%%%%%%%%%%%%%%
%---------------------------------------------------------------
\newcommand{\IntRange}[1]{\left\llbracket #1 \right\rrbracket}
\newcommand{\IRX}[1]{\left\llbracket #1 \right\rrbracket}%
\newcommand{\IRY}[2]{\left\llbracket #1:#2 \right\rrbracket}

%---------------------------------------------------------------
%%%%%%%%%%%%%%%%%%%%%%%%%%%%%%%%%%%%%%%%%%%%%%%%%%%%%%%%%%%%%%%%%%%%%%%

\newcommand{\obj}{\Mh{f}}%
\newcommand{\Wg}{\Mh{\omega}}%
\newcommand{\TotalW}{\Mh{W}}%

\newcommand{\rankY}[2]{{#1}_{ \permut{#2}}}%%
\newcommand{\areaX}[1]{\Mh{\mathrm{area}}\pth{#1}}%%

\newcommand{\mybinom}[2]%
{\Bigl(\begin{array}{@{}c@{}}#1\\#2\end{array}\Bigr)}

\providecommand{\Set}[2]{\left\{ #1 \;\middle\vert\; #2 \right\}}
\newcommand*{\biga}[1]{\vcenter{\hbox{\scalebox{1.5}{\ensuremath#1}}}}
\newcommand{\CNF}{\TermI{CNF}\xspace}%
\newcommand{\Tree}{\Mh{T}}%

\DeclareSymbolFont{sfletters}{OML}{cmbrm}{m}{it}

\DeclareMathSymbol{\salpha}{\mathord}{sfletters}{"0B}

%%%%%%%%%%%%%%%%%%%%%%%%%%%%%%%%%%%%%%%%%%%%%%%%%%%%%%%%%%%%%%%%%%%%

\providecommand{\ProblemInd}[2]{\textsf{\textcolor{red}{#1}\index{Problem!#2}}}
\newcommand{\NPProblemInd}[4]{
   \begin{center}
      \vspace{0.6cm}
       \fbox{
          \begin{minipage}{0.9\linewidth}
              \vspace{-0.6cm} \hspace{-0.35cm}
              {\large \bfseries \ProblemInd{#1}{#2}}
              \vspace{+0.2cm}%

              \noindent{\textcolor{blue}{\bfseries \textsf{Instance}}:} {#3}
              
              \noindent{\textcolor{blue}{\bfseries \textsf{Question}:}} {#4}
          \end{minipage}
       }
   \end{center}
   \smallskip
   }

\newcommand{\NPProblemOpt}[4]{\NPProblemInd{#1}{#2}{#3}{#4}}

%%%%%%%%%%%%%%%%%%%%%%%%%%%%%%%%%%%%%%%%%%%%%%%%%%%%%%%%%%%%%%%%%%%%

\newcommand{\Hastad}{H{\aa}stad\xspace}

\newcommand{\ANNbr}{\Term{ANN}\index{approximate nearest neighbor}\xspace}
\newcommand{\LSH}{\TermI{L{S}H}\xspace}

\newcommand{\DS}{\EuScript{D}} % DataStructure
\newcommand{\DSC}{\mathsf{D}} % DataStructure

\newcommand{\distH}[2]{\DistChar_{H}\pth{#1, #2}}
\newcommand{\distHY}[2]{\DistChar_{H}\pth{#1, #2}}

\newcommand{\DSAprxNearNbr}{\ensuremath{\DataStructure{{\approx%
         \text{Near{}}}}}%
   \index{near neighbor!data-structure!approximate}%
}
\newcommand{\DataStructure}[1]{\ensuremath{\DS_{#1}}}
\newcommand{\DistNN}[2]{\DistChar\pth{#1,#2}}
\newcommand{\pntc}{\Mh{{p}}}%

\newcommand{\seq}{{\Mh{M}}}%
\newcommand{\cTimes}{\Mh{\beta}}%

\newcommand{\DSTimes}{\Mh{L}}%
\newcommand{\rad}{\Mh{r}}%

\newcommand{\HypercubeX}[1]{\Hypercube^{#1}}
\newcommand{\kFunctions}[2]{\mathrm{combine}\pth{#1,#2}}
\newcommand{\WEqualChar}{\Bumpeq}
\newcommand{\WEqual}[2]{{#1}\WEqualChar{#2}}
\newcommand{\FamilyAmp}{\FamilyC_{\WEqualChar}}
\newcommand{\PSPACE}{\ComplexityClass{PSPACE}\xspace}%

\newcommand{\PS}{\Mh{P}}%

\newcommand{\PSetFar}{\ensuremath{\PntSet^{}_{\geq}\xspace}}%
\newcommand{\DA}{\Mh{D}}%
\newcommand{\uniqX}[1]{\Mh{\mathrm{uniq}}\pth{#1}}%
\newcommand{\diffX}[1]{\mathop{}\!\mathrm{d}#1}

\newcommand{\Frq}{\Mh{F}}%
\newcommand{\frq}{\Mh{f}}%
\newcommand{\STRM}{\Mh{\mathcal{S}}}%
\newcommand{\strm}{\Mh{{s}}}%

\providecommand{\Turan}{Tur\'an\xspace}
\newcommand{\fGrowX}[1]{\Mh{G}^{}_{\Dim}\pth{#1}}
\newcommand{\CompressFunc}{\mathtt{Compress}}
\newcommand{\RSA}{\Term{RSA}\xspace}
\newcommand{\MST}{\Term{MST}\xspace}
\providecommand{\CClassPP}{\CClass{P{}P}\xspace}
\providecommand{\ID}{\mathrm{id}}

\newcommand{\EncodeFunc}{\text{\ttfamily{Enc}}}
\newcommand{\DecodeFunc}{\text{\ttfamily{Dec}}}
\newcommand{\discX}[1]{\mathrm{disc}\pth{#1}}%

\newcommand{\propertyC}{\psi}%
\newcommand{\propertyX}[1]{\propertyC\pth{#1}}%
\newcommand{\numC}{\Gamma}%

\newcommand{\degreeX}[1]{\mathrm{degree}\pth{#1}}

\newcommand{\ind}{m}

\providecommand{\nequiv}{\not\equiv}
\providecommand{\sqcupplus}{\sqcup^+}

\newcommand{\AlgorithmI}[1]{{\protect\color{RedViolet}{#1}}}

\MakePerPage{footnote} % Resets footnote counter on every page

\allowdisplaybreaks % Allows page breaks inside multiline math environments

%-----------------------------------------------------------------------------
% Bibliography & Document Begin
%-----------------------------------------------------------------------------
\IfFileExists{\jobname.bib}{%
   \addbibresource{\jobname.bib}%
}{%
   \addbibresource{shortcuts.bib}%
   \addbibresource{geometry.bib}%
}
\makeindex
\begin{document}
%
\fi
\ChapterNumPage{17}{144}


\Chapter{Min Cut}

\QuotePExt{\textbf{To acknowledge the corn} - This purely American
   expression means to admit the losing of an argument, especially in
   regard to a detail; to retract; to admit defeat.  It is over a
   hundred years old.  Andrew Stewart, a member of Congress, is said
   to have mentioned it in a speech in 1828. He said that haystacks
   and cornfields were sent by Indiana, Ohio and Kentucky to
   Philadelphia and New York.  Charles A. Wickliffe, a member from
   Kentucky questioned the statement by commenting that haystacks and
   cornfields could not walk.  Stewart then pointed out that he did
   not mean literal haystacks and cornfields, but the horses, mules,
   and hogs for which the hay and corn were raised.  Wickliffe then
   rose to his feet, and said, ``Mr. Speaker, I acknowledge the
   corn''.}%
{Funk, Earle}{A Hog on Ice and Other Curious Expressions}


\noindent%
\ChapterPrefix{}%
\ChapterGen{}

~

% \bigskip%

\bigskip%


\section{Branching processes -- Galton-Watson Process}

% The problem that initiated the study of these processes goes back to
\subsection{The problem} %

In the 19\th{} century, Victorians were worried that aristocratic
surnames were disappearing, as family names passed on only through the
male children. As such, a family with no male children had its family
name disappear. So, imagine that the number of male children of an individual is
an independent random variable $X \in \brc{0,1,2,\ldots}$. Starting
with a single person, their family (as far as male children are
concerned) is a random tree with the number of children of a node being
distributed according to $X$. We continue recursively in constructing
this tree, again sampling the number of children for each current
leaf according to the distribution of $X$.  It is not hard to see that
a family disappears if $\Ex{X} \leq 1$, and it has a constant
probability of surviving if $\Ex{X} >1$.%
\remX{In branching process theory, extinction occurs almost surely (probability $1$) if $\Ex{X} < 1$, or if $\Ex{X} = 1$ with $\Prob{X=1} < 1$. In the degenerate critical case where $\Prob{X=1} = 1$, the population size remains $1$ forever and the family never disappears.}

Francis Galton asked the question of what is the probability of such a
blue-blood family name surviving, and this question was answered by
Henry William Watson~\cite{wg-pef-1875}. The Victorians were worried
about strange things; see~\cite{g-wwr-1869} for a provocatively titled
article from the period, and~\cite{s-wnrsf-12} for a more recent take
on this issue.

Of course, since infant mortality has decreased dramatically (as has the
number of aristocratic males dying to maintain the British empire), the
probability of a family name disappearing is now much lower than it was
in the 19\th{} century.  Interestingly, countries with family names that
were introduced a long time ago have very few surnames (e.g., Koreans
have only about 250 surnames, and three surnames form 45\% of the population). On
the other hand, countries that introduced surnames more recently have
dramatically more surnames (for example, the Dutch have surnames only
for the last 200 years, and there are $68,000$ different family
names).

Here we are going to look \chgY{on}{at} a very specific variant of this
problem. \chgY{Imagine that starting}{Imagine starting} with a single male.
\chgY{A male has exactly two children, and one of them is a male with probability half (i.e., the $Y$-chromosome is being passed only to its male children).}{Each male has exactly two children, each independently being male with probability $1/2$ (the $Y$-chromosome is passed only to male children).}
As such, the natural question is what is the probability that $h$
generations down, there is a male \chgY{decedent that all his ancestors}{descendant all of whose paternal ancestors} are male (i.e., \chgY{it caries}{he carries} the original family name, and the original
$Y$-chromosome).



\subsection{On coloring trees}
\seclab{min:cut:coloring:trees}

Let $T_h$ be a complete binary tree of height $h$. We randomly color
its edges black and white. Namely, for each edge we independently
choose its color to be either black or white, with equal probability
(say, black indicates the child is male). We are interested in the
event that there exists a path from the root of $T_h$ to one of its
\chgY{leafs, that is}{leaves that is} all black. Let $\Event_h$ denote this event, and let
$\rho_h = \Prob{\Event_h}$. Observe that $\rho_0 =1$ and $\rho_1 =
3/4$ (see below).

To bound this probability, consider the root $u$ of $T_h$ and its two
children $u_l$ and $u_r$. The probability that there is a black path
from $u_l$ to one of \chgY{its children}{the leaves of its subtree} is $\rho_{h-1}$, and as such, the
probability that there is a black path from $u$ through $u_l$ to a
leaf of the subtree of $u_l$ is $\Prob{\text{the edge } \MakeSBig u {}
   u_l \text{ is colored black}} \cdot \rho_{h-1} = \rho_{h-1}/2$. As
such, the probability that there is no black path through $u_l$ is
$1-\rho_{h-1}/2$. As such, the probability of not having a black path
from $u$ to a leaf (through \chgY{either children}{either child}) is
$(1-\rho_{h-1}/2)^2$. In particular, \chgY{there desired probability, is}{the desired probability is} the
complement; that is
\begin{equation*}
    \rho_h%
    =%
    1 - \pth{1-\frac{\rho_{h-1}}{2}}^2 %
    =%
    \frac{\rho_{h-1}}{2} \pth{2-\frac{\rho_{h-1}}{2}} = \rho_{h-1} -
    \frac{\rho_{h-1}^2}{4} = f \bigl( \rho_{h-1} \bigr)
    \qquad
    \text{ for }
    \qquad%
    f(x) = x -x^2/4.
\end{equation*}
The starting values are $\rho_0 = 1$, and $\rho_1 = 3/4$.

\begin{figure}[h]
    \centerline{%
       \IncludeGraphics[width=0.3\linewidth]%
       {\File{figs/stupid_function}}%
    }
    \caption{A graph of the function $f(x) = x -x^2/4$.}
\end{figure}

\begin{lemma}
    \lemlab{black:path}%
    %
    We have that $\rho_h \geq 1/(h+1)$.
\end{lemma}
\begin{proof}[Feel free to skip reading.]
    The proof is by induction. For $h=1$, we have
    $\rho_1 = 3/4 \geq 1/(1+1)$.
    
    Observe that $\rho_h = f\pth{\rho_{h-1} }$ for $f(x) = x -x^2/4$,
    and $f'(x) = 1-x/2$. As such, $f'(x) > 0$ for $x\in[0,1]$ and
    $f(x)$ is increasing in the range $[0,1]$. As such, by induction,
    we have that
    \begin{equation*}
        \ds \rho_h = f\pth[]{ \rho_{h-1}} \geq f\pth[]{
           \frac{1}{(h-1)+1}} = \frac{1}{h} - \frac{1}{4 h^2}.        
    \end{equation*}
    We need
    to prove that $\rho_h \geq 1/(h+1)$, which is implied by the above
    if
    \begin{equation*}
      \frac{1}{h} - \frac{1}{4 h^2} \geq \frac{1}{h+1}
      \;\;\;\Leftrightarrow \;\;\; 4h(h+1) - (h+1) \geq 4h^2
      \;\;\;\Leftrightarrow \;\;\; 4h^2 +4h - h- 1 \geq 4h^2
      \;\;\;\Leftrightarrow \;\;\; 3h \geq 1,
  \end{equation*}
    which trivially holds.
\end{proof}


\begin{lemma}
    \lemlab{black:path:2}%
    % 
    We have that $\rho_h =O( 1/h)$.
\end{lemma}


\begin{proof}[Feel free to skip reading.]
    We prove the claim for \chgY{infinite}{an infinite} number of values of $h$ -- the
    claim then follows for all $h$ by fiddling with the constants.
    The claim trivially holds for small values of $h$. For any
    $j > 0$, let $h_j$ be the minimal index such that
    $\rho_{h_j} \leq 1/2^j$.  It is easy to verify that
    $\rho_{h_j} \geq 1/2^{j+1}$.  We claim (mysteriously) that
    \begin{equation*}
        \ds h_{j+1} - h_j%
        \leq%
        \frac{\rho_{h_{j}} - \rho_{h_{j+1}}} { (\rho_{h_{j+1}} )^2/4
        }.%        
    \end{equation*}
    Indeed, $\rho_{k+1}$ is the number resulting from removing
    $\rho_k^2/4$ from $\rho_k$. Namely, the sequence
    $\rho_1, \rho_2, \ldots$ is a monotonically decreasing sequence of
    numbers in the interval $[0,1]$, where the gaps between
    consecutive numbers \chgY{decreases}{decrease}. In particular, to get from
    $\rho_{h_{j}}$ to $\rho_{h_{j+1}}$, the gaps used were of size at
    least $\Delta = \bigl(\rho_{h_{j+1}} \bigr)^2/4$, which means that
    there are \chgY{at least}{at most} $(\rho_{h_{j}} - \rho_{h_{j+1}}) / \Delta$
    numbers in the series between these two elements.%
    \remX{Corrected `at least' to `at most', and added the missing factor $1/4$ in $\Delta = (\rho_{h_{j+1}})^2/4$. If each step size is $\Delta_k \geq \Delta$, the number of steps cannot exceed $(\rho_{h_j} - \rho_{h_{j+1}})/\Delta$. Furthermore, since $h_j \leq 2^{j+6}$, for any $h$ with $h_{j-1} \leq h < h_j$, we have $h < h_j \leq 64 \cdot 2^j \implies 2^j \geq h/64$, so $\rho_h \leq \rho_{h_{j-1}} \leq 2^{-(j-1)} \leq 128/h$. This proves $\rho_h = O(1/h)$ for all $h \geq 1$ without needing separate constant-fiddling. Alternatively, a direct induction shows $\rho_h \leq 4/(h+4)$ for all $h \geq 1$, since $f(4/(h+4)) = \frac{4(h+3)}{(h+4)^2} < \frac{4}{h+5}$.}
    As such, since
    $\rho_{h_{j}} \leq 1/2^{j}$ and $\rho_{h_{j+1}} \geq 1/2^{j+2}$,
    we have
    \begin{equation*}
        h_{j+1} - h_j%
        \leq%
        \frac{\rho_{h_{j}} - \rho_{h_{j+1}}} {
           ( \rho_{h_{j+1}} )^2/4 }%
        \leq%
        \frac{1/2^{j} - 1/2^{j+2}} { 1/2^{2(j+2)+2}}%
        \leq%
        2^{2j + 6} /2^{j}
        =%
        2^{j+6}.
    \end{equation*}
    This implies that 
    \begin{math}
        h_j \leq (h_j - h_{j-1}) + (h_{j-1} - h_{j-2}) + \cdots + (h_1
        - h_0) \leq 2^{j+6}.
    \end{math}
    % Arguing similarly, we have
    % \begin{equation*}
    %     h_{j+2} - h_j%
    %     \geq%
    %     \frac{\rho_{h_{j}} - \rho_{h_{j+2}}} {
    %        (\rho_{h_{j}})^2/4 }%
    %     \geq%
    %     \frac{1/2^{j+1} - 1/2^{j+2}} { 1/2^{2j+2}}%
    %     =%
    %     2^{2j+2} \pth{1/2^{j+1} - 1/2^{j+2}}
    %     =%
    %     2^{j+1} - 2^{j}%
    %     =%
    %     2^j.
    % \end{equation*}
    % We conclude that
    % \begin{math}
    %     h_j \geq h_j - h_{j-2} \geq 2^{j-2}.
    % \end{math}
    As such, we have
    $\rho_{h_j} \leq 1/2^{j} \leq 2^6/2^{j+6} \leq 2^6/h_j$, which
    implies the claim.
\end{proof}



\section{Min Cut}

\subsection{Problem Definition}



\begin{figure}[t]
    \centering
    \IncGraphPage{figs/}{cut_example}{1}
    \caption{An example of a cut in a graph.}
    \figlab{cut:example}
\end{figure}

Let $\Graph=(\Vertices, \Edges)$ be an undirected graph with $n$
vertices and $m$ edges.  We are interested in \emphOnly{cuts}\index{cuts} in
$\Graph$.

\begin{defn}
    A \emphiC{cut} in $\Graph$ is a partition of the vertices of
    $\Vertices$ into two sets $S$ and $\Vertices \setminus S$,
    where the edges of the cut are
    \begin{equation*}
        (S, \Vertices \setminus S)%
        =%
        \brc{ uv \sep{ u \in S, v \in \Vertices \setminus S,
              \text{ and } uv \in \Edges}},
    \end{equation*}
    where $S \neq \emptyset$ and $\Vertices \setminus S \ne
    \emptyset$. We will refer to the number of edges in the cut
    $(S, \Vertices \setminus S)$ as the \emph{size of the cut}.
    For an example of a cut, see \figref{cut:example}.
\end{defn}


We are interested in the problem of computing the \emphic{minimum
   cut}{cut!minimum} (i.e., \emphOnly{mincut}\index{mincut}), that is, the cut in the
graph with minimum cardinality.  Specifically, we would like to find
the set $S\subseteq \Vertices$ such that $(S, \Vertices \setminus S)$
is as small as possible, and \chgY{$S$ is neither empty nor $\Vertices
\setminus S$ is empty}{neither $S$ nor $\Vertices \setminus S$ is empty}.



\subsection{Some Definitions}

We remind the reader of the following concepts.  The
\emphOnly{conditional probability}\index{conditional probability} of $X$ given $Y$ is $ \Prob{X=x\sep{
      Y=y}} = \nfrac{\Prob{(X=x)\cap(Y=y)}}{\Prob{Y=y}}$.  An
equivalent, useful restatement of this is that
\begin{equation}
    \Prob{\bigl. (X=x)\cap(Y=y)} = \Prob{X=x\sep{ Y=y \MakeSBig }}
    \cdot \Prob{Y=y}.%
    \eqlab{conditional:probability}
\end{equation}

% \smallskip

The following is easy to prove by induction using
\Eqref{conditional:probability}.
\begin{lemma}
    Let $\Event_{1},\ldots,\Event_{n}$ be $n$ events which are not
    necessarily independent. Then,
    \begin{align*}
      \Prob{\Bigl. \bigcap_{i=1}^{n}\Event_{i}} %
      =%
      \Prob{\bigl. \Event_{1}} \cdot \Prob{\Event_{2}\sep{ \Event_{1}}}
      \cdot \Prob{\Event_{3}\sep{\Event_{1}\cap\Event_{2}}} \cdot \ldots \cdot
      \Prob{\Event_{n} \sep{ \Event_{1}\cap\ldots\cap\Event_{n-1}}}.
    \end{align*}
\end{lemma}


\section{The Algorithm}

\begin{figure}[t]
   \centering
   \TwoFiguresBase{\IncGraphPage{figs/}{cut_example}{3}}{(a)}
                  {\IncGraphPage{figs/}{cut_example}{4}}{(b)}                  
   \caption{(a) A contraction of the edge $xy$. (b) The resulting graph.}
   \figlab{cut}
\end{figure}


The basic operation used by the algorithm is \emphic{edge
   contraction}{contraction!edge}, depicted in \figref{cut}.  We take
an edge $e=x y$ in $\Graph$ and merge the two vertices into a single
vertex.  The new resulting graph is denoted by $\Graph/ x y$.  Note
that we remove self-loops created by the contraction.  \chgY{However, since
the resulting graph is no longer a regular graph, it has parallel
edges -- namely, it is a multi-graph.  We represent a multi-graph, as
a regular graph with multiplicities on the edges.}{However, since
the resulting graph is no longer a simple graph, it may have parallel
edges---namely, it is a multigraph. We represent a multigraph as
a graph with multiplicities (or weights) on the edges.}%
\remX{Replaced `regular graph' with `simple graph'. In graph theory, a regular graph is one where every vertex has the same degree ($d$-regular), whereas a graph without multi-edges and self-loops is called a simple graph.}
See \figref{multi:graph}.


\begin{figure}[t]
   \centering
   \TwoFiguresBase{\IncGraphPage{figs/}{cut_example}{5}}{(a)}
                  {\IncGraphPage{figs/}{cut_example}{6}}{(b)}%
   \caption{(a) A multi-graph. (b) A minimum cut in the resulting multi-graph.}%
   \figlab{multi:graph}%
\end{figure}

The edge contraction operation can be implemented in $O(n)$ time for a
graph with $n$ vertices. This is done by merging the adjacency lists
of the two vertices being contracted, and then using hashing to do the
fix-ups (i.e., we need to fix the adjacency lists of the vertices that
are connected to the two vertices).

Note that the cut is now computed counting multiplicities (i.e., if
$e$ is in the cut and it has weight $w$, then the contribution of $e$
to the cut weight is $w$).


\begin{observation}
    A set of vertices in $\Graph/x y$ corresponds to a set of vertices
    in the graph $\Graph$. Thus a cut in $\Graph/x y$ always
    corresponds to a valid cut in $\Graph$. However, there are cuts in
    $\Graph$ that do not exist in $\Graph/ x y$. For example, the cut
    $S = \brc{x}$ does not exist in $\Graph/ x y$. As such, the size
    of the minimum cut in $\Graph/x y$ is at least as large as the
    \chgY{minimum cut in $\Graph$ (as long as $\Graph/x y$ has at least one
    edge). Since any cut in $\Graph/x y$ has a corresponding cut of
    the same cardinality in $\Graph$.}{minimum cut in $\Graph$ (as long as $\Graph/x y$ has at least one
    edge), since any cut in $\Graph/x y$ has a corresponding cut of
    the same cardinality in $\Graph$.}
\end{observation}

Our algorithm works by repeatedly performing edge contractions. This
is beneficial as this shrinks the underlying graph, and we \chgY{would}{can}
compute the cut in the resulting (smaller) graph. An ``extreme''
example of this is shown in \figref{contraction:cut}, where we
contract the graph into a single edge, which (in turn) corresponds to
a cut in the original graph. (It might help the reader to think about
each vertex in the contracted graph as corresponding to a connected
component in the original graph.)
\begin{figure}[p]
    \begin{center}
        \begin{tabular}{*{4}{c}}
            \IncGraphPage{figs}{cut_example}{2}
            &
            \IncGraphPage{figs}{cut_example}{3}
            &
            \IncGraphPage{figs}{cut_example}{7}
            &
            \IncGraphPage{figs}{cut_example}{8}\\
            (a) & (b) & (c) & (d)\\
            \IncGraphPage{figs}{cut_example}{9}
            &
            \IncGraphPage{figs}{cut_example}{10}
            &
            \IncGraphPage{figs}{cut_example}{11}
            &
            \IncGraphPage{figs}{cut_example}{12}\\
            (e) & (f) & (g) & (h)\\
            &\IncGraphPage{figs}{cut_example}{13}
            &
            \IncGraphPage{figs}{cut_example}{14}\\
            &(i) & (j)
        \end{tabular}
    \end{center}
    \caption{(a) Original graph. \chgY{(b)--(j) a sequence of contractions
       in the graph, and (h) the cut in the original graph,
       corresponding to the single edge in (h). Note that the cut of
       (h) is not a mincut in the original graph.}{(b)--(i) a sequence of contractions
       in the graph resulting in the two-vertex multigraph in (i), and (j) the cut in the original graph
       corresponding to the single multi-edge in (i). Note that the cut in (j)
       is not a mincut in the original graph.}}%
    \remX{Corrected subfigure references: (b)--(i) show the contractions leading to the 2-vertex multigraph in (i) (page 13), and (j) shows the corresponding cut in the original graph (page 14). The original caption mistakenly referenced (h) for both.}
    \figlab{contraction:cut}
\end{figure}

\figref{contraction:cut} also demonstrates the problem with taking
this approach. Indeed, the resulting cut is not the minimum cut in the
graph.

So, why did the algorithm fail to find the minimum cut in this
case?\footnote{Naturally, if the algorithm had succeeded in finding
   the minimum cut, this would have been \textbf{our} success.}  The
failure occurs because of the contraction at \figref{contraction:cut}
(e), as we had contracted an edge in the minimum cut. In the new
graph, depicted in \figref{contraction:cut} (f), there is no longer a
cut of size 3, and all cuts are of size 4 or more. Specifically, the
algorithm succeeds only if it does not contract an edge in the minimum
cut.

\begin{observation}
    Let $e_{1},\ldots,e_{n-2}$ be a sequence of edges in $\Graph$,
    such that none of them is in the minimum cut, and such that
    $\Graph'=\Graph/\brc{ e_{1},\ldots,e_{n-2}} $ is a single
    multi-edge. Then, this multi-edge corresponds to a minimum cut in
    $\Graph$.
    
    \obslab{contract}
\end{observation}

%\parpic[r]{%
%   \begin{minipage}{0.5\linewidth}%
\begin{figure}[p]
    \centerline{%
    \begin{algorithmEnv}
           {\bfseries{Algorithm }} \MinCut{}($\Graph$) \+ \\
           $\Graph_{0}\leftarrow \Graph$\\
           $i=0$\\
           \While{} $\Graph_{i}$ has more than two vertices \Do{}\+\\
           Pick randomly an edge $e_{i}$ from the edges of
           $\Graph_{i}$\\
           $\Graph_{i+1}\leftarrow \Graph_{i}/e_{i}$\\
           $i\leftarrow i+1$ \- \\
           Let $(S,\Vertices \setminus S)$ be the cut in the original graph\\
           ~~~~~~~~~~~~~ corresponding
           to the multi-edge in $\Graph_{i}$\\
           \Return $(S, \Vertices \setminus S)$.%
       \end{algorithmEnv}%
       }%
       % \vspace*{-0.4cm}%
       
       \captionof{figure}{The minimum cut algorithm.}
       \figlab{mincut:algorithm}%
       % \end{minipage}%
\end{figure}
%}

Note that the claim in the above observation is only in one
direction. We might be able to still compute a minimum cut even if we
contract an edge in a minimum cut, the reason being that a minimum cut
is not unique. In particular, another minimum cut \chgY{might survived}{might survive} the
sequence of contractions that destroyed other minimum cuts.


Using \obsref{contract} in an algorithm is problematic, since the
argumentation is circular: how can we find a sequence of edges that
are not in the cut without knowing what the cut is?  The way to slice
the Gordian knot here is to randomly select an edge at each stage,
and contract this random edge.


See \figref{mincut:algorithm} for the resulting algorithm \MinCut{}.


\subsection{Analysis}

\subsubsection{The probability of success}
Naturally, if we are extremely lucky, the algorithm would never pick
an edge in the mincut, and the algorithm would succeed. The ultimate
question here is what is the probability of success. If it is
relatively ``large'' then this algorithm is useful since we can run it
several times, and return the best result computed. If on the other
hand, this probability is tiny, then we are working in vain since this
approach would not work.

\begin{lemma}
    If a graph $\Graph$ has a minimum cut of size $k$ and $\Graph$ has
    $n$ vertices, then $|\Edges(\Graph)|\geq\frac{kn}{2}$.
\end{lemma}

\begin{proof}
    Each vertex degree is at least $k$, otherwise the vertex itself
    would form a cut of size smaller than $k$. As such, there
    are at least $\sum_{v \in \Vertices} \degreeX{v} / 2 \geq n
    k /2$ edges in the graph.
\end{proof}

\begin{lemma}
    \lemlab{random:edge}%
    %
    Fix a specific minimum cut $C = (S,\overline{S})$ in the graph.
    If we pick an edge $e$ uniformly at random from a graph $\Graph$, then with probability at most $\nfrac{2}{n}$ it belongs
    to the minimum cut $C$.
\end{lemma}

\begin{proof}
    There are at least $n k/2$ edges in the graph and exactly $k$
    edges in the minimum cut. Thus, the probability of picking an edge
    from the minimum cut is at most $k/(n k/2)=2/n$.
\end{proof}


% \medskip

\providecommand\Ccancel[2][black]{%
   \renewcommand\CancelColor{\color{#1}}\cancel{#2}}

The following lemma shows (surprisingly) that \MinCut{} succeeds with
reasonable probability.


\begin{lemma}
    \lemlab{mincut:success}%
    %
    \MinCut{} outputs the mincut with probability $\displaystyle \geq
    \frac{2}{n ( n-1 ) }$.    
\end{lemma}
\begin{proof}
    Fix a target minimum cut $C$.
    Let $\Event_{i}$ be the event that $e_{i}$ is not in $C$.%
    \remX{Removed a stray digit `7' that appeared immediately before `\textbackslash begin\{proof\}' in the original source, and clarified that $\Event_i$ is the event that $e_i \notin C$ for a fixed target minimum cut $C$. Conditioned on $\Event_0 \cap \dots \cap \Event_{i-1}$, the cut $C$ survives intact in $\Graph_i$, so the mincut in $\Graph_i$ is still of size $k$, ensuring $|\Edges(\Graph_i)| \geq k(n-i)/2$.}
    By \obsref{contract}, \MinCut{} outputs the
    minimum cut if the events $\Event_{0},\ldots,\Event_{n-3}$ all
    happen (namely, all edges picked are outside the minimum cut).
    
    By \lemref{random:edge}, it holds
    \begin{math}
        \displaystyle \Prob{\Event_{i}\sep{ \Event_{0} \cap
              \Event_{1}\cap\ldots\cap\Event_{i-1}}} \geq
        1-\frac{2}{\cardin{\Vertices(\Graph_{i})}} = 1-\frac{2}{n-i}.
    \end{math}
    Implying that
    \begin{align*}
        \Delta%
        &=%
        \Prob{\Event_{0}\cap\ldots\cap\Event_{n-3}}%
        % \\&%
        =%
        \Prob{\Event_{0}} \cdot \Prob{\Event_{1}\sep{\Event_{0}}}
        \cdot \Prob{\Event_{2}\sep{\Event_{0}\cap\Event_{1}}} \cdot
        \ldots \cdot \Prob{\Event_{n-3}\sep{
              \Event_{0}\cap\ldots\cap\Event_{n-4}}}.
    \end{align*}
    As such, we have
    \begin{align*}
      \Delta%
      &\geq %
        \prod_{i=0}^{n-3}\pth{1-\frac{2}{n-i}}%
        =%
        \prod_{i=0}^{n-3}\frac{n-i-2}{n-i}%
      \\&%
      =%
      \frac{\CCancel[green]{n-2}}{n}
      \cdot \frac{\Ccancel{n-3}}{n-1}
      \cdot \frac{\Ccancel[orange]{n-4}}{\CCancel[green]{n-2}}%
      \cdot
      \frac{\bcancel{n-5}}{\cancel{n-3}}
      \cdot
      \frac{\Ccancel[red]{n-6}}{\Ccancel[orange]{n-4}}
      \cdot
      \cdots \cdot \frac{\Ccancel[purple]{3}}{\Ccancel[yellow]{5}} \cdot
      \frac{2}{\CCancel[brown]{4}}
      \cdot \frac{1}{\CCancel[purple]{3}}%
      \\&
      =%
      \frac{2}{n (n-1)}.
    \end{align*}
\end{proof}

\subsubsection{Running time analysis.}

\begin{observation}
    \MinCut{} runs in $O( n^{2} )$ time.
\end{observation}

\begin{observation}
    The algorithm always outputs a cut, and the cut is not smaller
    than the minimum cut.
\end{observation}

\begin{defn}\index{Probability!Amplification}
    (informal) Amplification is the process of running an experiment
    again and again till the things we want to happen, with good
    probability, do happen.
\end{defn}

Let \MinCutRep{} be the algorithm that runs \MinCut{} $n(n-1)$ times
and return the minimum cut computed in all those independent
executions of \MinCut{}.

\begin{lemma}
    The probability that \MinCutRep{} fails to return the minimum cut
    is $< 0.14$.
    
    \lemlab{mincut:amplification}
\end{lemma}
\begin{proof}
    The probability of failure of \MinCut{} to output the mincut in each
    execution is at most $1-\frac{2}{n( n-1)}$, by
    \lemref{mincut:success}. Now, \MinCutRep{} fails only if all the $n
    (n-1)$ executions of \MinCut{} fail. But these executions are
    independent, as such, the probability of this happening is at most
    \begin{align*}
        \pth{1-\frac{2}{n
              (n-1)}}^{n(n-1)}\leq\exp\pth{-\frac{2}{n(n-1)}\cdot
           n(n-1)}=\exp(-2)<0.14,
    \end{align*}
    since $1-x\leq e^{-x}$ for $0\leq x\leq 1$.
\end{proof}

\begin{theorem}   
    One can compute the minimum cut in $O(n^{4})$ time with constant
    probability to get a correct result. In $O\pth{n^{4}\log n}$ time
    the minimum cut is returned with high probability.
\end{theorem}

%%%%%%%%%%%%%%%%%%%%%%%%%%%%%%%%%%%%%%%%%%%%%%%%%%%%%%%%%%%%%%%%%%%%%%%

\subsection{An alternative implementation using MST}

\paragraph*{The algorithm.}
The above algorithm can be restated as follows. Randomly assign
weights to the edges of $\Graph$ (say, by picking numbers in
$[0,1]$). Next, compute the \MST{} $\Tree$ of the graph according to
these weights. Remove the heaviest edge in the \MST{}. The resulting
partition of $\Tree$ into two trees corresponds to a cut in the
original graph. Return this cut as a candidate to be the minimum cut. 

\paragraph*{The analysis.}
To see that this algorithm is equivalent to \MinCut{}
(\figref{mincut:algorithm}), observe that the contraction algorithm
simulates Kruskal's \MST{} algorithm when run on randomly weighted
edges. First, imagine implementing \MinCut{} so that it keeps parallel
edges. Then, the edges connecting two vertices that are not
contracted are exactly the edges between the two connected
components. Picking a random edge to contract is equivalent to
picking the edge with the minimum random weight. Thus, the \MST{}
algorithm here just simulates \MinCut{} (or vice versa). 

\paragraph*{A small optimization.}
It is possible to compute the heaviest edge in the \MST{}, and the
partition it induces in (deterministic) linear time -- it is a 
nice example of the search and prune technique.  
\begin{exercise}
    Given a graph $\Graph$ with weights on the edges, show how to compute
    the maximum weight edge in the \MST{} of $\Graph$ in $O(n + m)$ time,
    where $n$ and $m$ are the number of vertices and edges of $\Graph$,
    respectively.%
    \remX{Exercise 3.9 asks for finding the maximum-weight edge in the MST in $O(n+m)$ deterministic time without computing the full MST. This is the bottleneck spanning tree problem, solved in $O(n+m)$ time by Camerini's prune-and-search algorithm (Camerini, 1978) using linear-time median finding to contract or discard half the edges at each step. Adding a hint pointing to bottleneck spanning trees / prune-and-search would clarify the intended approach.}
\end{exercise}

Thus, this yields an $O(n + m)$ implementation of \MinCut{}. We get the
following result.

\begin{lemma}
    \lemlab{mincut:MST}%
    %
    \MinCut{} can be implemented to run in $O(n + m)$ time, and it
    outputs the mincut with probability
    $\displaystyle \geq \frac{2}{n ( n-1 ) }$.
\end{lemma}


%%%%%%%%%%%%%%%%%%%%%%%%%%%%%%%%%%%%%%%%%%%%%%%%%%%%%%%%%%%%%%%%%%%%%%%
%%%%%%%%%%%%%%%%%%%%%%%%%%%%%%%%%%%%%%%%%%%%%%%%%%%%%%%%%%%%%%%%%%%%%%%
%%%%%%%%%%%%%%%%%%%%%%%%%%%%%%%%%%%%%%%%%%%%%%%%%%%%%%%%%%%%%%%%%%%%%%%

\section{A faster algorithm}

The algorithm presented in the previous section is extremely simple.
Which raises the question of whether we can get a faster
algorithm\footnote{This would require a more involved algorithm, that
   is life.}?

So, why does \MinCutRep{} need so many executions? Well, the probability
of success in the first $\nu$ iterations is
\begin{align}
  \Prob{\Event_{0}\cap\ldots\cap\Event_{\nu-1}}
  &\geq 
    \prod_{i=0}^{\nu-1}\pth{1-\frac{2}{n-i}}
    =  \prod_{i=0}^{\nu-1}\frac{n-i-2}{n-i} \nonumber\\
  &=  \frac{n-2}{n}\cdot \frac{n-3}{n-1}\cdot \frac{n-4}{n-2}\cdots =
    \frac{(n-\nu)(n-\nu-1)}{n\cdot(n-1)}.  \eqlab{first}
\end{align}
Namely, this probability deteriorates very quickly toward the end of
the execution, when the graph becomes small enough. (To see this,
observe that for $\nu=n/2$, the probability of success is roughly
$1/4$, but for $\nu=n-\sqrt{n}$ the probability of success is roughly
$1/n$.)


So, the key observation is that as the graph gets smaller the
probability of making a bad choice increases.  So, instead of doing the
amplification from the outside of the algorithm, we will run the new
algorithm more times when the graph is smaller. Namely, we put the
amplification directly into the algorithm.


\begin{figure}
    \hfill
    \begin{algorithmEnv}
        \!\!\Contract( $\Graph, t$ )\\
        \Procbegin{}  \+  \\
        \While{} $\cardin{\Vertices(\Graph)}>t$ \Do\+\\%
        \remX{Corrected `\textbackslash cardin\{(G)\}' to `\textbackslash cardin\{\textbackslash Vertices(\textbackslash Graph)\}'. The original typeset as $|(G)|>t$.}%
        Pick a random edge $e$ in $\Graph$.\\
        $\Graph\leftarrow \Graph/e$\-\\
        \Return{} $\Graph$\- \\
        \End{}
    \end{algorithmEnv}
    \hfill
    \begin{algorithmEnv}
        \noindent%
        \!\!\FastCut{}($\Graph=(\Vertices,\Edges)$) \+\\
        $\Graph$ -- multi-graph \- \\
        \Procbegin{}\+\\
        $n\leftarrow\cardin{\Vertices(\Graph)}$\\
        \If{} $n\leq 6$ \Then\+ \\
        Compute (via brute force) minimum cut\\
        of $\Graph$ and return
        cut.\-\\
        $t\leftarrow \ceil{ 1+ n/\sqrt{2}}$\\
        $H_{1} \leftarrow \Contract{}(\Graph,t)$\\
        $H_{2} \leftarrow \Contract{}(\Graph,t)$\\
        /* \Contract{} is randomized!!! {*}/\\
        $X_{1}\leftarrow \FastCut{}\pth{H_{1}}$,\\
        $X_{2}\leftarrow \FastCut{}\pth{H_{2}}$\\
        \Return minimum cut out of $X_{1}$ and
        $X_{2}$.\- \\
        \End{}
    \end{algorithmEnv}
    \hfill
    
    \caption{\Contract{}($ \Graph, t$) shrinks $\Graph$ till it has only
       $t$ vertices. \FastCut{} computes the minimum cut using
       \Contract{}.}
    \figlab{contract}
\end{figure}

The basic new operation we use is \Contract{}, depicted in
\figref{contract}, which also depicts the new algorithm \FastCut{}.

% \begin{minipage}[c]{0.45\textwidth}

% \end{minipage}



\begin{lemma}
    The running time of \FastCut{}$(\Graph)$ is $O\pth{n^{2}\log n}$, where
    $n=\cardin{\Vertices(\Graph)}$.
    
    \lemlab{fast:cut:running:time}
\end{lemma}
\begin{proof}
    Well, we perform two calls to \Contract{}$(\Graph,t)$ which takes
    $O(n^{2})$ time. And then we perform two recursive calls on the
    resulting graphs.  We have
    \begin{equation*}
        T(n)=O(n^{2}) + 2 T\bigl(  n/ \sqrt{2} \ts \bigr).
    \end{equation*}
    The solution to this recurrence is $O\pth{n^{2}\log n}$ as one can
    easily (and should) verify.
\end{proof}



\begin{exercise}
    Show that one can modify \FastCut{} so that it uses only
    $O(n^{2})$ space.
\end{exercise}

\begin{lemma}
    The probability that
    \Contract{}$\bigl( \Graph, \ceil{1 + n/\sqrt{2}} \ts \bigr)$ does not
    contract the minimum cut is at least $1/2$.
    
    Namely, the probability that the minimum cut in the contracted
    graph is still a minimum cut in the original graph is at least
    $1/2$.
    
    \lemlab{easy}
\end{lemma}
\begin{proof}
    Just plug in $\nu= n -t = n - \ceil{1 + n/\sqrt{2} }$ into
    \Eqref{first}. We have
    \begin{align*}
        \Prob{\Bigl. \Event_{0}\cap\ldots\cap\Event_{n-t-1}}\geq &
        \frac{t(t-1)}{n\cdot(n-1)} = \frac{\ceil{1 + n/\sqrt{2} }
           \pth{ \ceil{1 + n/\sqrt{2} } - 1}}{n(n-1)} \geq
        \frac{1}{2}.
    \end{align*}%
    \remX{Corrected the upper index in the intersection from $\Event_{n-t}$ to $\Event_{n-t-1}$. Since $\nu = n-t$ edge contractions are performed starting with index $0$, the sequence of events is $\Event_0, \dots, \Event_{\nu-1} = \Event_{n-t-1}$.}
\end{proof}

The following lemma bounds the probability of success. 
%A more elegant
%argument is given in \secref{min:cut:coloring:trees} below.

\begin{lemma}
    \lemlab{fast:cut:probability}
    % 
    \FastCut{} finds the minimum cut with probability $\Omega\pth[]{1/\log n}$.    
\end{lemma}
\begin{proof}
    Let $T_h$ be the recursion tree of the algorithm of depth $h =
    \Theta( \log n)$.  Color an edge of the recursion tree black if the
    contraction does not contract any edge of the target minimum cut $C$.%
    \remX{Clarified `contraction succeeded' to explicitly state that the contraction preserves the target minimum cut $C$.}
    Clearly, the algorithm succeeds if there is
    a path from the root to a leaf that is all black. This is exactly
    the settings of \lemref{black:path}, and we conclude that the
    probability of success is at least $1/(h+1) = \Theta(1/\log n)$,
    as desired.
\end{proof}

\begin{exercise}
    Prove that running \FastCut{} repeatedly $c\cdot\log^{2}n$ times
    guarantees that the algorithm outputs the minimum cut with
    probability $\geq1-1/n^{2}$ for a sufficiently large constant $c$.
\end{exercise} %

\begin{theorem}
    One can compute the minimum cut in a graph $\Graph$ with $n$
    vertices in $O(n^2 \log^3 n)$ time. The algorithm succeeds with
    probability $\geq 1-1/n^2$.
\end{theorem}
\begin{proof}
    We do amplification on \FastCut{} by running it $O(\log^2 n)$
    times. The running time bound follows from
    \lemref{fast:cut:running:time}. The bound on the probability
    follows from \lemref{fast:cut:probability}, and using the
    amplification analysis as done in \lemref{mincut:amplification}
    for \MinCutRep{}.
\end{proof}






\section{Bibliographical Notes}

The \MinCut{} algorithm was developed by David Karger during his PhD
thesis at Stanford. The fast algorithm is joint work with Clifford
Stein.  The basic algorithm of the mincut is described in \cite[pages~7--9]{mr-ra-95}, the faster algorithm is described in \cite[pages~289--295]{mr-ra-95}.

\paragraph*{Galton-Watson process.}
The idea of using coloring of the edges of a tree to analyze \FastCut{}
might be new (i.e., \secref{min:cut:coloring:trees}).
% It is inspired by \emphOnly{Galton-Watson processes}\index{Galton-Watson processes} (which is a
% special case of a branching process).


% \bibliographystyle{alpha} \bibliography{shortcuts,geometry}

% -----------

% 
% \end{document}
\ChapterEnd{}

% -----------
